%=====================================================================
%  Direction Is Not Enough: Trigger Salience and Adaptive Collusion
%  against Byzantine-Robust Federated Intrusion Detection
%
%  Manuscript prepared for the Journal of Information Security and
%  Applications (Elsevier).
%
%  Build:  make          (or: latexmk -pdf main.tex)
%  Gate :  make gate     (refuses to pass while any \PENDING remains)
%
%  Layout: [3p] is single column at journal density, matching the
%  companion Computers & Security submission. For the two-column look
%  JISA uses in print, swap the class line for
%      \documentclass[final,5p,times,twocolumn]{elsarticle}
%  and see README.txt for the float and figure widths that go with it.
%=====================================================================
\documentclass[3p]{elsarticle}

% --- maths ----------------------------------------------------------
\usepackage{amsmath,amssymb,amsfonts}

% --- floats, figures, tables ----------------------------------------
\usepackage{graphicx}
\usepackage{booktabs}
\usepackage{multirow}

% --- typography -----------------------------------------------------
\usepackage[T1]{fontenc}
\usepackage{lmodern}       % bold small caps for \textsc{GradCos} in bold contexts
\usepackage{textcomp}
\usepackage{xcolor}
\usepackage{microtype}
\usepackage{url}
\usepackage{xurl}          % let the long repository URL break across lines
\usepackage[hidelinks]{hyperref}

% NB: no \usepackage{cite} -- elsarticle loads natbib itself, and the two
% clash silently, renumbering or dropping citations.

% Uncomment for a review submission if the editor asks for line numbers:
% \usepackage{lineno}\linenumbers

\graphicspath{{figs/}}

% Marks a result whose experiment is specified and implemented but has not yet
% been run to completion. NOTHING carrying this may go to a venue. Remove the
% macro entirely once results/remediation.log reports the pipeline complete.
\newcommand{\PENDING}{\textcolor{red}{\textbf{[PENDING RUN]}}}

\newcommand{\repourl}{https://github.com/usamamustafa36/federated-backdoor-gradcos}

% Every number quoted in this paper is defined here, generated from the committed
% result files by paper/make_numbers.py. Do not type a result into the prose or a
% table cell: add it to the generator instead. A cell that has not been run yet
% expands to \PENDING, so an unfinished claim is loud rather than silently wrong.
%   python3 paper/aggregate_multiseed.py && python3 -m experiments.make_tables
%   python3 paper/make_numbers.py --strict     # submission gate: exits 1 if any \PENDING
% GENERATED FILE — DO NOT EDIT. Regenerate with: python3 paper/make_numbers.py
% Source: results/summary_multiseed.csv, results/summary.csv
% Any \PENDING below is a cell that has not been run yet.

% ---- multi-seed headline (CICIoT2023, IID, m=0.2, p=0.1, seeds 42-46) --
\newcommand{\msCleanCAfedavg}{0.702}  % fedavg clean accuracy
\newcommand{\msCleanFonefedavg}{0.561}  % fedavg clean macro-F1
\newcommand{\msBSRfedavgFixed}{0.968}
\newcommand{\msBSRsdfedavgFixed}{0.012}
\newcommand{\msDRfedavgFixed}{\PENDING{}}
\newcommand{\msDRsdfedavgFixed}{\PENDING{}}
\newcommand{\msFPRfedavgFixed}{\PENDING{}}
\newcommand{\msFPRsdfedavgFixed}{\PENDING{}}
\newcommand{\msCAfedavgFixed}{0.700}
\newcommand{\msCAsdfedavgFixed}{0.013}
\newcommand{\msFPRpctfedavgFixed}{\PENDING{}}
\newcommand{\msBSRminfedavgFixed}{0.955}
\newcommand{\msBSRmaxfedavgFixed}{0.987}
\newcommand{\msSalfedavgFixed}{\PENDING{}}  % measured salience gap
\newcommand{\msNseedsfedavgFixed}{5}
\newcommand{\msBSRfedavgStatistical}{1.000}
\newcommand{\msBSRsdfedavgStatistical}{0.000}
\newcommand{\msDRfedavgStatistical}{\PENDING{}}
\newcommand{\msDRsdfedavgStatistical}{\PENDING{}}
\newcommand{\msFPRfedavgStatistical}{\PENDING{}}
\newcommand{\msFPRsdfedavgStatistical}{\PENDING{}}
\newcommand{\msCAfedavgStatistical}{0.715}
\newcommand{\msCAsdfedavgStatistical}{0.018}
\newcommand{\msFPRpctfedavgStatistical}{\PENDING{}}
\newcommand{\msBSRminfedavgStatistical}{1.000}
\newcommand{\msBSRmaxfedavgStatistical}{1.000}
\newcommand{\msSalfedavgStatistical}{\PENDING{}}  % measured salience gap
\newcommand{\msNseedsfedavgStatistical}{5}
\newcommand{\msBSRfedavgTemporal}{1.000}
\newcommand{\msBSRsdfedavgTemporal}{0.000}
\newcommand{\msDRfedavgTemporal}{\PENDING{}}
\newcommand{\msDRsdfedavgTemporal}{\PENDING{}}
\newcommand{\msFPRfedavgTemporal}{\PENDING{}}
\newcommand{\msFPRsdfedavgTemporal}{\PENDING{}}
\newcommand{\msCAfedavgTemporal}{0.698}
\newcommand{\msCAsdfedavgTemporal}{0.021}
\newcommand{\msFPRpctfedavgTemporal}{\PENDING{}}
\newcommand{\msBSRminfedavgTemporal}{1.000}
\newcommand{\msBSRmaxfedavgTemporal}{1.000}
\newcommand{\msSalfedavgTemporal}{\PENDING{}}  % measured salience gap
\newcommand{\msNseedsfedavgTemporal}{5}
\newcommand{\msCleanCAkrum}{0.703}  % krum clean accuracy
\newcommand{\msCleanFonekrum}{0.560}  % krum clean macro-F1
\newcommand{\msBSRkrumFixed}{0.948}
\newcommand{\msBSRsdkrumFixed}{0.016}
\newcommand{\msDRkrumFixed}{0.18}
\newcommand{\msDRsdkrumFixed}{0.05}
\newcommand{\msFPRkrumFixed}{0.20}
\newcommand{\msFPRsdkrumFixed}{0.01}
\newcommand{\msCAkrumFixed}{0.693}
\newcommand{\msCAsdkrumFixed}{0.009}
\newcommand{\msFPRpctkrumFixed}{20}
\newcommand{\msBSRminkrumFixed}{0.927}
\newcommand{\msBSRmaxkrumFixed}{0.970}
\newcommand{\msSalkrumFixed}{\PENDING{}}  % measured salience gap
\newcommand{\msNseedskrumFixed}{5}
\newcommand{\msBSRkrumStatistical}{1.000}
\newcommand{\msBSRsdkrumStatistical}{0.000}
\newcommand{\msDRkrumStatistical}{0.24}
\newcommand{\msDRsdkrumStatistical}{0.05}
\newcommand{\msFPRkrumStatistical}{0.19}
\newcommand{\msFPRsdkrumStatistical}{0.01}
\newcommand{\msCAkrumStatistical}{0.697}
\newcommand{\msCAsdkrumStatistical}{0.011}
\newcommand{\msFPRpctkrumStatistical}{19}
\newcommand{\msBSRminkrumStatistical}{0.999}
\newcommand{\msBSRmaxkrumStatistical}{1.000}
\newcommand{\msSalkrumStatistical}{\PENDING{}}  % measured salience gap
\newcommand{\msNseedskrumStatistical}{5}
\newcommand{\msBSRkrumTemporal}{1.000}
\newcommand{\msBSRsdkrumTemporal}{0.000}
\newcommand{\msDRkrumTemporal}{0.20}
\newcommand{\msDRsdkrumTemporal}{0.05}
\newcommand{\msFPRkrumTemporal}{0.20}
\newcommand{\msFPRsdkrumTemporal}{0.01}
\newcommand{\msCAkrumTemporal}{0.698}
\newcommand{\msCAsdkrumTemporal}{0.013}
\newcommand{\msFPRpctkrumTemporal}{20}
\newcommand{\msBSRminkrumTemporal}{1.000}
\newcommand{\msBSRmaxkrumTemporal}{1.000}
\newcommand{\msSalkrumTemporal}{\PENDING{}}  % measured salience gap
\newcommand{\msNseedskrumTemporal}{5}
\newcommand{\msCleanCAflame}{0.499}  % flame clean accuracy
\newcommand{\msCleanFoneflame}{0.334}  % flame clean macro-F1
\newcommand{\msBSRflameFixed}{0.010}
\newcommand{\msBSRsdflameFixed}{0.003}
\newcommand{\msDRflameFixed}{0.64}
\newcommand{\msDRsdflameFixed}{0.14}
\newcommand{\msFPRflameFixed}{0.40}
\newcommand{\msFPRsdflameFixed}{0.04}
\newcommand{\msCAflameFixed}{0.538}
\newcommand{\msCAsdflameFixed}{0.106}
\newcommand{\msFPRpctflameFixed}{40}
\newcommand{\msBSRminflameFixed}{0.007}
\newcommand{\msBSRmaxflameFixed}{0.015}
\newcommand{\msSalflameFixed}{\PENDING{}}  % measured salience gap
\newcommand{\msNseedsflameFixed}{5}
\newcommand{\msBSRflameStatistical}{0.863}
\newcommand{\msBSRsdflameStatistical}{0.137}
\newcommand{\msDRflameStatistical}{0.42}
\newcommand{\msDRsdflameStatistical}{0.08}
\newcommand{\msFPRflameStatistical}{0.46}
\newcommand{\msFPRsdflameStatistical}{0.02}
\newcommand{\msCAflameStatistical}{0.464}
\newcommand{\msCAsdflameStatistical}{0.068}
\newcommand{\msFPRpctflameStatistical}{46}
\newcommand{\msBSRminflameStatistical}{0.636}
\newcommand{\msBSRmaxflameStatistical}{0.974}
\newcommand{\msSalflameStatistical}{\PENDING{}}  % measured salience gap
\newcommand{\msNseedsflameStatistical}{5}
\newcommand{\msBSRflameTemporal}{0.863}
\newcommand{\msBSRsdflameTemporal}{0.225}
\newcommand{\msDRflameTemporal}{0.46}
\newcommand{\msDRsdflameTemporal}{0.07}
\newcommand{\msFPRflameTemporal}{0.45}
\newcommand{\msFPRsdflameTemporal}{0.02}
\newcommand{\msCAflameTemporal}{0.556}
\newcommand{\msCAsdflameTemporal}{0.040}
\newcommand{\msFPRpctflameTemporal}{45}
\newcommand{\msBSRminflameTemporal}{0.464}
\newcommand{\msBSRmaxflameTemporal}{0.990}
\newcommand{\msSalflameTemporal}{\PENDING{}}  % measured salience gap
\newcommand{\msNseedsflameTemporal}{5}
\newcommand{\msCleanCAfltrust}{0.556}  % fltrust clean accuracy
\newcommand{\msCleanFonefltrust}{0.426}  % fltrust clean macro-F1
\newcommand{\msBSRfltrustFixed}{0.004}
\newcommand{\msBSRsdfltrustFixed}{0.000}
\newcommand{\msDRfltrustFixed}{0.00}
\newcommand{\msDRsdfltrustFixed}{0.00}
\newcommand{\msFPRfltrustFixed}{0.00}
\newcommand{\msFPRsdfltrustFixed}{0.00}
\newcommand{\msCAfltrustFixed}{0.556}
\newcommand{\msCAsdfltrustFixed}{0.054}
\newcommand{\msFPRpctfltrustFixed}{0}
\newcommand{\msBSRminfltrustFixed}{0.004}
\newcommand{\msBSRmaxfltrustFixed}{0.004}
\newcommand{\msSalfltrustFixed}{\PENDING{}}  % measured salience gap
\newcommand{\msNseedsfltrustFixed}{5}
\newcommand{\msBSRfltrustStatistical}{0.007}
\newcommand{\msBSRsdfltrustStatistical}{0.010}
\newcommand{\msDRfltrustStatistical}{0.00}
\newcommand{\msDRsdfltrustStatistical}{0.00}
\newcommand{\msFPRfltrustStatistical}{0.00}
\newcommand{\msFPRsdfltrustStatistical}{0.00}
\newcommand{\msCAfltrustStatistical}{0.577}
\newcommand{\msCAsdfltrustStatistical}{0.067}
\newcommand{\msFPRpctfltrustStatistical}{0}
\newcommand{\msBSRminfltrustStatistical}{0.000}
\newcommand{\msBSRmaxfltrustStatistical}{0.022}
\newcommand{\msSalfltrustStatistical}{\PENDING{}}  % measured salience gap
\newcommand{\msNseedsfltrustStatistical}{5}
\newcommand{\msBSRfltrustTemporal}{0.782}
\newcommand{\msBSRsdfltrustTemporal}{0.147}
\newcommand{\msDRfltrustTemporal}{0.00}
\newcommand{\msDRsdfltrustTemporal}{0.00}
\newcommand{\msFPRfltrustTemporal}{0.00}
\newcommand{\msFPRsdfltrustTemporal}{0.00}
\newcommand{\msCAfltrustTemporal}{0.538}
\newcommand{\msCAsdfltrustTemporal}{0.066}
\newcommand{\msFPRpctfltrustTemporal}{0}
\newcommand{\msBSRminfltrustTemporal}{0.562}
\newcommand{\msBSRmaxfltrustTemporal}{0.892}
\newcommand{\msSalfltrustTemporal}{\PENDING{}}  % measured salience gap
\newcommand{\msNseedsfltrustTemporal}{5}
\newcommand{\msCleanCAfoolsgold}{0.703}  % foolsgold clean accuracy
\newcommand{\msCleanFonefoolsgold}{0.560}  % foolsgold clean macro-F1
\newcommand{\msBSRfoolsgoldFixed}{0.369}
\newcommand{\msBSRsdfoolsgoldFixed}{0.489}
\newcommand{\msDRfoolsgoldFixed}{0.45}
\newcommand{\msDRsdfoolsgoldFixed}{0.27}
\newcommand{\msFPRfoolsgoldFixed}{0.00}
\newcommand{\msFPRsdfoolsgoldFixed}{0.00}
\newcommand{\msCAfoolsgoldFixed}{0.701}
\newcommand{\msCAsdfoolsgoldFixed}{0.016}
\newcommand{\msFPRpctfoolsgoldFixed}{0}
\newcommand{\msBSRminfoolsgoldFixed}{0.003}
\newcommand{\msBSRmaxfoolsgoldFixed}{0.926}
\newcommand{\msSalfoolsgoldFixed}{\PENDING{}}  % measured salience gap
\newcommand{\msNseedsfoolsgoldFixed}{5}
\newcommand{\msBSRfoolsgoldStatistical}{1.000}
\newcommand{\msBSRsdfoolsgoldStatistical}{0.001}
\newcommand{\msDRfoolsgoldStatistical}{0.00}
\newcommand{\msDRsdfoolsgoldStatistical}{0.00}
\newcommand{\msFPRfoolsgoldStatistical}{0.00}
\newcommand{\msFPRsdfoolsgoldStatistical}{0.00}
\newcommand{\msCAfoolsgoldStatistical}{0.713}
\newcommand{\msCAsdfoolsgoldStatistical}{0.024}
\newcommand{\msFPRpctfoolsgoldStatistical}{0}
\newcommand{\msBSRminfoolsgoldStatistical}{0.999}
\newcommand{\msBSRmaxfoolsgoldStatistical}{1.000}
\newcommand{\msSalfoolsgoldStatistical}{\PENDING{}}  % measured salience gap
\newcommand{\msNseedsfoolsgoldStatistical}{5}
\newcommand{\msBSRfoolsgoldTemporal}{1.000}
\newcommand{\msBSRsdfoolsgoldTemporal}{0.000}
\newcommand{\msDRfoolsgoldTemporal}{0.00}
\newcommand{\msDRsdfoolsgoldTemporal}{0.00}
\newcommand{\msFPRfoolsgoldTemporal}{0.00}
\newcommand{\msFPRsdfoolsgoldTemporal}{0.00}
\newcommand{\msCAfoolsgoldTemporal}{0.699}
\newcommand{\msCAsdfoolsgoldTemporal}{0.020}
\newcommand{\msFPRpctfoolsgoldTemporal}{0}
\newcommand{\msBSRminfoolsgoldTemporal}{1.000}
\newcommand{\msBSRmaxfoolsgoldTemporal}{1.000}
\newcommand{\msSalfoolsgoldTemporal}{\PENDING{}}  % measured salience gap
\newcommand{\msNseedsfoolsgoldTemporal}{5}
\newcommand{\msCleanCAgradcos}{0.705}  % gradcos clean accuracy
\newcommand{\msCleanFonegradcos}{0.562}  % gradcos clean macro-F1
\newcommand{\msBSRgradcosFixed}{0.004}
\newcommand{\msBSRsdgradcosFixed}{0.003}
\newcommand{\msDRgradcosFixed}{0.97}
\newcommand{\msDRsdgradcosFixed}{0.01}
\newcommand{\msFPRgradcosFixed}{0.01}
\newcommand{\msFPRsdgradcosFixed}{0.01}
\newcommand{\msCAgradcosFixed}{0.700}
\newcommand{\msCAsdgradcosFixed}{0.017}
\newcommand{\msFPRpctgradcosFixed}{0}
\newcommand{\msBSRmingradcosFixed}{0.001}
\newcommand{\msBSRmaxgradcosFixed}{0.007}
\newcommand{\msSalgradcosFixed}{\PENDING{}}  % measured salience gap
\newcommand{\msNseedsgradcosFixed}{5}
\newcommand{\msBSRgradcosStatistical}{1.000}
\newcommand{\msBSRsdgradcosStatistical}{0.000}
\newcommand{\msDRgradcosStatistical}{0.11}
\newcommand{\msDRsdgradcosStatistical}{0.05}
\newcommand{\msFPRgradcosStatistical}{0.12}
\newcommand{\msFPRsdgradcosStatistical}{0.03}
\newcommand{\msCAgradcosStatistical}{0.707}
\newcommand{\msCAsdgradcosStatistical}{0.015}
\newcommand{\msFPRpctgradcosStatistical}{12}
\newcommand{\msBSRmingradcosStatistical}{1.000}
\newcommand{\msBSRmaxgradcosStatistical}{1.000}
\newcommand{\msSalgradcosStatistical}{\PENDING{}}  % measured salience gap
\newcommand{\msNseedsgradcosStatistical}{5}
\newcommand{\msBSRgradcosTemporal}{1.000}
\newcommand{\msBSRsdgradcosTemporal}{0.000}
\newcommand{\msDRgradcosTemporal}{0.14}
\newcommand{\msDRsdgradcosTemporal}{0.05}
\newcommand{\msFPRgradcosTemporal}{0.12}
\newcommand{\msFPRsdgradcosTemporal}{0.03}
\newcommand{\msCAgradcosTemporal}{0.706}
\newcommand{\msCAsdgradcosTemporal}{0.022}
\newcommand{\msFPRpctgradcosTemporal}{12}
\newcommand{\msBSRmingradcosTemporal}{1.000}
\newcommand{\msBSRmaxgradcosTemporal}{1.000}
\newcommand{\msSalgradcosTemporal}{\PENDING{}}  % measured salience gap
\newcommand{\msNseedsgradcosTemporal}{5}
\newcommand{\msCAadvGradcos}{0.15}  % gradcos clean CA over fltrust
\newcommand{\msCAdeficitFlame}{0.20}  % flame clean CA below fedavg

% ---- adaptive attacker: static vs decorrelated vs adaptive -------------
\newcommand{\adBSRfedavgSimple}{0.968}
\newcommand{\adDRfedavgSimple}{\PENDING{}}
\newcommand{\adFPRfedavgSimple}{\PENDING{}}
\newcommand{\adCAfedavgSimple}{0.700}
\newcommand{\adMalCosfedavgSimple}{0.999}
\newcommand{\adHonCosfedavgSimple}{0.999}
\newcommand{\adNormRatiofedavgSimple}{1.00}
\newcommand{\adBSRfedavgDecorrelated}{0.957}
\newcommand{\adDRfedavgDecorrelated}{\PENDING{}}
\newcommand{\adFPRfedavgDecorrelated}{\PENDING{}}
\newcommand{\adCAfedavgDecorrelated}{0.690}
\newcommand{\adMalCosfedavgDecorrelated}{0.428}
\newcommand{\adHonCosfedavgDecorrelated}{0.999}
\newcommand{\adNormRatiofedavgDecorrelated}{1.32}
\newcommand{\adBSRfedavgAdaptive}{0.957}
\newcommand{\adDRfedavgAdaptive}{\PENDING{}}
\newcommand{\adFPRfedavgAdaptive}{\PENDING{}}
\newcommand{\adCAfedavgAdaptive}{0.697}
\newcommand{\adMalCosfedavgAdaptive}{0.420}
\newcommand{\adHonCosfedavgAdaptive}{0.999}
\newcommand{\adNormRatiofedavgAdaptive}{0.99}
\newcommand{\evBSRfedavgQuarter}{\PENDING{}}
\newcommand{\evMalCosfedavgQuarter}{\PENDING{}}
\newcommand{\evBSRfedavgHalf}{\PENDING{}}
\newcommand{\evMalCosfedavgHalf}{\PENDING{}}
\newcommand{\evBSRfedavgOne}{0.957}
\newcommand{\evMalCosfedavgOne}{0.420}
\newcommand{\evBSRfedavgTwo}{\PENDING{}}
\newcommand{\evMalCosfedavgTwo}{\PENDING{}}
\newcommand{\adBSRkrumSimple}{0.948}
\newcommand{\adDRkrumSimple}{0.18}
\newcommand{\adFPRkrumSimple}{0.20}
\newcommand{\adCAkrumSimple}{0.693}
\newcommand{\adMalCoskrumSimple}{\PENDING{}}
\newcommand{\adHonCoskrumSimple}{\PENDING{}}
\newcommand{\adNormRatiokrumSimple}{\PENDING{}}
\newcommand{\adBSRkrumDecorrelated}{0.001}
\newcommand{\adDRkrumDecorrelated}{1.00}
\newcommand{\adFPRkrumDecorrelated}{0.00}
\newcommand{\adCAkrumDecorrelated}{0.692}
\newcommand{\adMalCoskrumDecorrelated}{0.428}
\newcommand{\adHonCoskrumDecorrelated}{0.999}
\newcommand{\adNormRatiokrumDecorrelated}{1.32}
\newcommand{\adBSRkrumAdaptive}{0.001}
\newcommand{\adDRkrumAdaptive}{1.00}
\newcommand{\adFPRkrumAdaptive}{0.00}
\newcommand{\adCAkrumAdaptive}{0.692}
\newcommand{\adMalCoskrumAdaptive}{0.421}
\newcommand{\adHonCoskrumAdaptive}{0.999}
\newcommand{\adNormRatiokrumAdaptive}{0.99}
\newcommand{\evBSRkrumQuarter}{\PENDING{}}
\newcommand{\evMalCoskrumQuarter}{\PENDING{}}
\newcommand{\evBSRkrumHalf}{\PENDING{}}
\newcommand{\evMalCoskrumHalf}{\PENDING{}}
\newcommand{\evBSRkrumOne}{0.001}
\newcommand{\evMalCoskrumOne}{0.421}
\newcommand{\evBSRkrumTwo}{\PENDING{}}
\newcommand{\evMalCoskrumTwo}{\PENDING{}}
\newcommand{\adBSRflameSimple}{0.010}
\newcommand{\adDRflameSimple}{0.64}
\newcommand{\adFPRflameSimple}{0.40}
\newcommand{\adCAflameSimple}{0.538}
\newcommand{\adMalCosflameSimple}{\PENDING{}}
\newcommand{\adHonCosflameSimple}{\PENDING{}}
\newcommand{\adNormRatioflameSimple}{\PENDING{}}
\newcommand{\adBSRflameDecorrelated}{\PENDING{}}
\newcommand{\adDRflameDecorrelated}{\PENDING{}}
\newcommand{\adFPRflameDecorrelated}{\PENDING{}}
\newcommand{\adCAflameDecorrelated}{\PENDING{}}
\newcommand{\adMalCosflameDecorrelated}{\PENDING{}}
\newcommand{\adHonCosflameDecorrelated}{\PENDING{}}
\newcommand{\adNormRatioflameDecorrelated}{\PENDING{}}
\newcommand{\adBSRflameAdaptive}{\PENDING{}}
\newcommand{\adDRflameAdaptive}{\PENDING{}}
\newcommand{\adFPRflameAdaptive}{\PENDING{}}
\newcommand{\adCAflameAdaptive}{\PENDING{}}
\newcommand{\adMalCosflameAdaptive}{\PENDING{}}
\newcommand{\adHonCosflameAdaptive}{\PENDING{}}
\newcommand{\adNormRatioflameAdaptive}{\PENDING{}}
\newcommand{\evBSRflameQuarter}{\PENDING{}}
\newcommand{\evMalCosflameQuarter}{\PENDING{}}
\newcommand{\evBSRflameHalf}{\PENDING{}}
\newcommand{\evMalCosflameHalf}{\PENDING{}}
\newcommand{\evBSRflameOne}{\PENDING{}}
\newcommand{\evMalCosflameOne}{\PENDING{}}
\newcommand{\evBSRflameTwo}{\PENDING{}}
\newcommand{\evMalCosflameTwo}{\PENDING{}}
\newcommand{\adBSRfltrustSimple}{0.004}
\newcommand{\adDRfltrustSimple}{0.00}
\newcommand{\adFPRfltrustSimple}{0.00}
\newcommand{\adCAfltrustSimple}{0.556}
\newcommand{\adMalCosfltrustSimple}{\PENDING{}}
\newcommand{\adHonCosfltrustSimple}{\PENDING{}}
\newcommand{\adNormRatiofltrustSimple}{\PENDING{}}
\newcommand{\adBSRfltrustDecorrelated}{\PENDING{}}
\newcommand{\adDRfltrustDecorrelated}{\PENDING{}}
\newcommand{\adFPRfltrustDecorrelated}{\PENDING{}}
\newcommand{\adCAfltrustDecorrelated}{\PENDING{}}
\newcommand{\adMalCosfltrustDecorrelated}{\PENDING{}}
\newcommand{\adHonCosfltrustDecorrelated}{\PENDING{}}
\newcommand{\adNormRatiofltrustDecorrelated}{\PENDING{}}
\newcommand{\adBSRfltrustAdaptive}{\PENDING{}}
\newcommand{\adDRfltrustAdaptive}{\PENDING{}}
\newcommand{\adFPRfltrustAdaptive}{\PENDING{}}
\newcommand{\adCAfltrustAdaptive}{\PENDING{}}
\newcommand{\adMalCosfltrustAdaptive}{\PENDING{}}
\newcommand{\adHonCosfltrustAdaptive}{\PENDING{}}
\newcommand{\adNormRatiofltrustAdaptive}{\PENDING{}}
\newcommand{\evBSRfltrustQuarter}{\PENDING{}}
\newcommand{\evMalCosfltrustQuarter}{\PENDING{}}
\newcommand{\evBSRfltrustHalf}{\PENDING{}}
\newcommand{\evMalCosfltrustHalf}{\PENDING{}}
\newcommand{\evBSRfltrustOne}{\PENDING{}}
\newcommand{\evMalCosfltrustOne}{\PENDING{}}
\newcommand{\evBSRfltrustTwo}{\PENDING{}}
\newcommand{\evMalCosfltrustTwo}{\PENDING{}}
\newcommand{\adBSRfoolsgoldSimple}{0.369}
\newcommand{\adDRfoolsgoldSimple}{0.45}
\newcommand{\adFPRfoolsgoldSimple}{0.00}
\newcommand{\adCAfoolsgoldSimple}{0.701}
\newcommand{\adMalCosfoolsgoldSimple}{\PENDING{}}
\newcommand{\adHonCosfoolsgoldSimple}{\PENDING{}}
\newcommand{\adNormRatiofoolsgoldSimple}{\PENDING{}}
\newcommand{\adBSRfoolsgoldDecorrelated}{\PENDING{}}
\newcommand{\adDRfoolsgoldDecorrelated}{\PENDING{}}
\newcommand{\adFPRfoolsgoldDecorrelated}{\PENDING{}}
\newcommand{\adCAfoolsgoldDecorrelated}{\PENDING{}}
\newcommand{\adMalCosfoolsgoldDecorrelated}{\PENDING{}}
\newcommand{\adHonCosfoolsgoldDecorrelated}{\PENDING{}}
\newcommand{\adNormRatiofoolsgoldDecorrelated}{\PENDING{}}
\newcommand{\adBSRfoolsgoldAdaptive}{\PENDING{}}
\newcommand{\adDRfoolsgoldAdaptive}{\PENDING{}}
\newcommand{\adFPRfoolsgoldAdaptive}{\PENDING{}}
\newcommand{\adCAfoolsgoldAdaptive}{\PENDING{}}
\newcommand{\adMalCosfoolsgoldAdaptive}{\PENDING{}}
\newcommand{\adHonCosfoolsgoldAdaptive}{\PENDING{}}
\newcommand{\adNormRatiofoolsgoldAdaptive}{\PENDING{}}
\newcommand{\evBSRfoolsgoldQuarter}{\PENDING{}}
\newcommand{\evMalCosfoolsgoldQuarter}{\PENDING{}}
\newcommand{\evBSRfoolsgoldHalf}{\PENDING{}}
\newcommand{\evMalCosfoolsgoldHalf}{\PENDING{}}
\newcommand{\evBSRfoolsgoldOne}{\PENDING{}}
\newcommand{\evMalCosfoolsgoldOne}{\PENDING{}}
\newcommand{\evBSRfoolsgoldTwo}{\PENDING{}}
\newcommand{\evMalCosfoolsgoldTwo}{\PENDING{}}
\newcommand{\adBSRgradcosSimple}{0.004}
\newcommand{\adDRgradcosSimple}{0.97}
\newcommand{\adFPRgradcosSimple}{0.01}
\newcommand{\adCAgradcosSimple}{0.700}
\newcommand{\adMalCosgradcosSimple}{\PENDING{}}
\newcommand{\adHonCosgradcosSimple}{\PENDING{}}
\newcommand{\adNormRatiogradcosSimple}{\PENDING{}}
\newcommand{\adBSRgradcosDecorrelated}{\PENDING{}}
\newcommand{\adDRgradcosDecorrelated}{\PENDING{}}
\newcommand{\adFPRgradcosDecorrelated}{\PENDING{}}
\newcommand{\adCAgradcosDecorrelated}{\PENDING{}}
\newcommand{\adMalCosgradcosDecorrelated}{\PENDING{}}
\newcommand{\adHonCosgradcosDecorrelated}{\PENDING{}}
\newcommand{\adNormRatiogradcosDecorrelated}{\PENDING{}}
\newcommand{\adBSRgradcosAdaptive}{\PENDING{}}
\newcommand{\adDRgradcosAdaptive}{\PENDING{}}
\newcommand{\adFPRgradcosAdaptive}{\PENDING{}}
\newcommand{\adCAgradcosAdaptive}{\PENDING{}}
\newcommand{\adMalCosgradcosAdaptive}{\PENDING{}}
\newcommand{\adHonCosgradcosAdaptive}{\PENDING{}}
\newcommand{\adNormRatiogradcosAdaptive}{\PENDING{}}
\newcommand{\evBSRgradcosQuarter}{\PENDING{}}
\newcommand{\evMalCosgradcosQuarter}{\PENDING{}}
\newcommand{\evBSRgradcosHalf}{\PENDING{}}
\newcommand{\evMalCosgradcosHalf}{\PENDING{}}
\newcommand{\evBSRgradcosOne}{\PENDING{}}
\newcommand{\evMalCosgradcosOne}{\PENDING{}}
\newcommand{\evBSRgradcosTwo}{\PENDING{}}
\newcommand{\evMalCosgradcosTwo}{\PENDING{}}

% ---- GradCos residual variant, static vs adaptive ----------------------
\newcommand{\resBSRSimpleFixed}{0.003}
\newcommand{\resBSRsdSimpleFixed}{0.001}
\newcommand{\resDRSimpleFixed}{0.99}
\newcommand{\resFPRSimpleFixed}{0.03}
\newcommand{\resCASimpleFixed}{0.695}
\newcommand{\resBSRSimpleStatistical}{0.085}
\newcommand{\resBSRsdSimpleStatistical}{0.131}
\newcommand{\resDRSimpleStatistical}{0.99}
\newcommand{\resFPRSimpleStatistical}{0.03}
\newcommand{\resCASimpleStatistical}{0.704}
\newcommand{\resBSRSimpleTemporal}{0.367}
\newcommand{\resBSRsdSimpleTemporal}{0.170}
\newcommand{\resDRSimpleTemporal}{0.99}
\newcommand{\resFPRSimpleTemporal}{0.05}
\newcommand{\resCASimpleTemporal}{0.698}
\newcommand{\resBSRAdaptiveFixed}{\PENDING{}}
\newcommand{\resBSRsdAdaptiveFixed}{\PENDING{}}
\newcommand{\resDRAdaptiveFixed}{\PENDING{}}
\newcommand{\resFPRAdaptiveFixed}{\PENDING{}}
\newcommand{\resCAAdaptiveFixed}{\PENDING{}}
\newcommand{\resBSRAdaptiveStatistical}{\PENDING{}}
\newcommand{\resBSRsdAdaptiveStatistical}{\PENDING{}}
\newcommand{\resDRAdaptiveStatistical}{\PENDING{}}
\newcommand{\resFPRAdaptiveStatistical}{\PENDING{}}
\newcommand{\resCAAdaptiveStatistical}{\PENDING{}}
\newcommand{\resBSRAdaptiveTemporal}{\PENDING{}}
\newcommand{\resBSRsdAdaptiveTemporal}{\PENDING{}}
\newcommand{\resDRAdaptiveTemporal}{\PENDING{}}
\newcommand{\resFPRAdaptiveTemporal}{\PENDING{}}
\newcommand{\resCAAdaptiveTemporal}{\PENDING{}}

% ---- magnitude and order-statistic rules (the non-cosine control family) 
\newcommand{\mgCleanCAmedian}{\PENDING{}}
\newcommand{\mgBSRmedianFixed}{\PENDING{}}
\newcommand{\mgCAmedianFixed}{\PENDING{}}
\newcommand{\mgSalmedianFixed}{\PENDING{}}
\newcommand{\mgBSRmedianStatistical}{\PENDING{}}
\newcommand{\mgCAmedianStatistical}{\PENDING{}}
\newcommand{\mgSalmedianStatistical}{\PENDING{}}
\newcommand{\mgBSRmedianTemporal}{\PENDING{}}
\newcommand{\mgCAmedianTemporal}{\PENDING{}}
\newcommand{\mgSalmedianTemporal}{\PENDING{}}
\newcommand{\mgCleanCAtrimmedmean}{\PENDING{}}
\newcommand{\mgBSRtrimmedmeanFixed}{\PENDING{}}
\newcommand{\mgCAtrimmedmeanFixed}{\PENDING{}}
\newcommand{\mgSaltrimmedmeanFixed}{\PENDING{}}
\newcommand{\mgBSRtrimmedmeanStatistical}{\PENDING{}}
\newcommand{\mgCAtrimmedmeanStatistical}{\PENDING{}}
\newcommand{\mgSaltrimmedmeanStatistical}{\PENDING{}}
\newcommand{\mgBSRtrimmedmeanTemporal}{\PENDING{}}
\newcommand{\mgCAtrimmedmeanTemporal}{\PENDING{}}
\newcommand{\mgSaltrimmedmeanTemporal}{\PENDING{}}
\newcommand{\mgCleanCArfa}{\PENDING{}}
\newcommand{\mgBSRrfaFixed}{\PENDING{}}
\newcommand{\mgCArfaFixed}{\PENDING{}}
\newcommand{\mgSalrfaFixed}{\PENDING{}}
\newcommand{\mgBSRrfaStatistical}{\PENDING{}}
\newcommand{\mgCArfaStatistical}{\PENDING{}}
\newcommand{\mgSalrfaStatistical}{\PENDING{}}
\newcommand{\mgBSRrfaTemporal}{\PENDING{}}
\newcommand{\mgCArfaTemporal}{\PENDING{}}
\newcommand{\mgSalrfaTemporal}{\PENDING{}}
\newcommand{\mgCleanCAbulyan}{\PENDING{}}
\newcommand{\mgBSRbulyanFixed}{\PENDING{}}
\newcommand{\mgCAbulyanFixed}{\PENDING{}}
\newcommand{\mgSalbulyanFixed}{\PENDING{}}
\newcommand{\mgBSRbulyanStatistical}{\PENDING{}}
\newcommand{\mgCAbulyanStatistical}{\PENDING{}}
\newcommand{\mgSalbulyanStatistical}{\PENDING{}}
\newcommand{\mgBSRbulyanTemporal}{\PENDING{}}
\newcommand{\mgCAbulyanTemporal}{\PENDING{}}
\newcommand{\mgSalbulyanTemporal}{\PENDING{}}
\newcommand{\mgCleanCAnormclip}{\PENDING{}}
\newcommand{\mgBSRnormclipFixed}{\PENDING{}}
\newcommand{\mgCAnormclipFixed}{\PENDING{}}
\newcommand{\mgSalnormclipFixed}{\PENDING{}}
\newcommand{\mgBSRnormclipStatistical}{\PENDING{}}
\newcommand{\mgCAnormclipStatistical}{\PENDING{}}
\newcommand{\mgSalnormclipStatistical}{\PENDING{}}
\newcommand{\mgBSRnormclipTemporal}{\PENDING{}}
\newcommand{\mgCAnormclipTemporal}{\PENDING{}}
\newcommand{\mgSalnormclipTemporal}{\PENDING{}}

% ---- trigger realizability audit (CICIoT2023) --------------------------
\newcommand{\rzSigFixedQtr}{0.50}
\newcommand{\rzDMFixedQtr}{11.59}
\newcommand{\rzNNFixedQtr}{1.79}
\newcommand{\rzSigFixedOne}{2.00}
\newcommand{\rzDMFixedOne}{29.09}
\newcommand{\rzNNFixedOne}{3.75}
\newcommand{\rzSigFixedThree}{6.00}
\newcommand{\rzDMFixedThree}{76.72}
\newcommand{\rzNNFixedThree}{10.68}
\newcommand{\rzSigFixedDonor}{11.12}
\newcommand{\rzDMFixedDonor}{2.62}
\newcommand{\rzNNFixedDonor}{6.77}
\newcommand{\rzPnnFixed}{5.78}  % benign 99th-pctile joint Mahalanobis
\newcommand{\rzSigStatisticalQtr}{0.50}
\newcommand{\rzDMStatisticalQtr}{0.75}
\newcommand{\rzNNStatisticalQtr}{1.01}
\newcommand{\rzSigStatisticalOne}{2.00}
\newcommand{\rzDMStatisticalOne}{1.12}
\newcommand{\rzNNStatisticalOne}{3.01}
\newcommand{\rzSigStatisticalThree}{6.00}
\newcommand{\rzDMStatisticalThree}{5.23}
\newcommand{\rzNNStatisticalThree}{8.50}
\newcommand{\rzSigStatisticalDonor}{15.83}
\newcommand{\rzDMStatisticalDonor}{2.95}
\newcommand{\rzNNStatisticalDonor}{15.25}
\newcommand{\rzPnnStatistical}{7.16}  % benign 99th-pctile joint Mahalanobis
\newcommand{\rzSigTemporalQtr}{0.50}
\newcommand{\rzDMTemporalQtr}{922}
\newcommand{\rzNNTemporalQtr}{0.52}
\newcommand{\rzSigTemporalOne}{2.00}
\newcommand{\rzDMTemporalOne}{1167}
\newcommand{\rzNNTemporalOne}{2.62}
\newcommand{\rzSigTemporalThree}{6.00}
\newcommand{\rzDMTemporalThree}{1826}
\newcommand{\rzNNTemporalThree}{8.44}
\newcommand{\rzSigTemporalDonor}{8.71}
\newcommand{\rzDMTemporalDonor}{2.56}
\newcommand{\rzNNTemporalDonor}{8.49}
\newcommand{\rzPnnTemporal}{3.96}  % benign 99th-pctile joint Mahalanobis
\newcommand{\rzUnstamped}{0.002}  % nearest-real-neighbour distance for a real attack flow
\newcommand{\rzExceedFixedQtr}{5}  % benign rows further from the benign mean than the 0.25x fixed trigger
\newcommand{\rzRefRows}{20{,}000}  % benign reference rows in the audit
\newcommand{\rzPctileFixedQtr}{99.97}  % where the 0.5-sigma fixed trigger sits in benign joint variation

% ---- breadth grid: mean BSR over the 9 attack cells of each trigger ----
\newcommand{\gridBSRfedavgFixed}{0.959}  % mean of 9 cells
\newcommand{\gridNfedavgFixed}{9}
\newcommand{\gridBSRfedavgStatistical}{1.000}  % mean of 9 cells
\newcommand{\gridNfedavgStatistical}{9}
\newcommand{\gridBSRfedavgTemporal}{1.000}  % mean of 9 cells
\newcommand{\gridNfedavgTemporal}{9}
\newcommand{\gridBSRkrumFixed}{0.559}  % mean of 9 cells
\newcommand{\gridNkrumFixed}{9}
\newcommand{\gridBSRkrumStatistical}{0.799}  % mean of 9 cells
\newcommand{\gridNkrumStatistical}{9}
\newcommand{\gridBSRkrumTemporal}{0.671}  % mean of 9 cells
\newcommand{\gridNkrumTemporal}{9}
\newcommand{\gridBSRflameFixed}{0.048}  % mean of 9 cells
\newcommand{\gridNflameFixed}{9}
\newcommand{\gridBSRflameStatistical}{0.906}  % mean of 9 cells
\newcommand{\gridNflameStatistical}{9}
\newcommand{\gridBSRflameTemporal}{0.705}  % mean of 9 cells
\newcommand{\gridNflameTemporal}{9}
\newcommand{\gridBSRfltrustFixed}{0.004}  % mean of 9 cells
\newcommand{\gridNfltrustFixed}{9}
\newcommand{\gridBSRfltrustStatistical}{0.001}  % mean of 9 cells
\newcommand{\gridNfltrustStatistical}{9}
\newcommand{\gridBSRfltrustTemporal}{0.568}  % mean of 9 cells
\newcommand{\gridNfltrustTemporal}{9}
\newcommand{\gridBSRfoolsgoldFixed}{0.026}  % mean of 1 cells
\newcommand{\gridNfoolsgoldFixed}{1}
\newcommand{\gridBSRfoolsgoldStatistical}{1.000}  % mean of 1 cells
\newcommand{\gridNfoolsgoldStatistical}{1}
\newcommand{\gridBSRfoolsgoldTemporal}{1.000}  % mean of 1 cells
\newcommand{\gridNfoolsgoldTemporal}{1}
\newcommand{\gridBSRgradcosFixed}{0.171}  % mean of 9 cells
\newcommand{\gridNgradcosFixed}{9}
\newcommand{\gridBSRgradcosStatistical}{0.893}  % mean of 9 cells
\newcommand{\gridNgradcosStatistical}{9}
\newcommand{\gridBSRgradcosTemporal}{0.999}  % mean of 9 cells
\newcommand{\gridNgradcosTemporal}{9}
\newcommand{\gridAttackCells}{27}  % undefended attack cells in the breadth grid
\newcommand{\gridMinBSRfedavg}{0.899}  % weakest undefended cell

% ---- malicious-ratio sweep (fixed trigger, mean over poison rates) -----
\newcommand{\ratBSRfedavgTen}{0.91}  % mean of 3 poison rates
\newcommand{\ratDRfedavgTen}{\PENDING{}}
\newcommand{\ratNfedavgTen}{3}
\newcommand{\ratBSRfedavgTwenty}{0.98}  % mean of 3 poison rates
\newcommand{\ratDRfedavgTwenty}{\PENDING{}}
\newcommand{\ratNfedavgTwenty}{3}
\newcommand{\ratBSRfedavgThirty}{0.99}  % mean of 3 poison rates
\newcommand{\ratDRfedavgThirty}{\PENDING{}}
\newcommand{\ratNfedavgThirty}{3}
\newcommand{\ratBSRkrumTen}{0.00}  % mean of 3 poison rates
\newcommand{\ratDRkrumTen}{1.00}
\newcommand{\ratNkrumTen}{3}
\newcommand{\ratBSRkrumTwenty}{0.69}  % mean of 3 poison rates
\newcommand{\ratDRkrumTwenty}{0.30}
\newcommand{\ratNkrumTwenty}{3}
\newcommand{\ratBSRkrumThirty}{0.98}  % mean of 3 poison rates
\newcommand{\ratDRkrumThirty}{0.13}
\newcommand{\ratNkrumThirty}{3}
\newcommand{\ratBSRflameTen}{0.01}  % mean of 3 poison rates
\newcommand{\ratDRflameTen}{0.65}
\newcommand{\ratNflameTen}{3}
\newcommand{\ratBSRflameTwenty}{0.02}  % mean of 3 poison rates
\newcommand{\ratDRflameTwenty}{0.49}
\newcommand{\ratNflameTwenty}{3}
\newcommand{\ratBSRflameThirty}{0.12}  % mean of 3 poison rates
\newcommand{\ratDRflameThirty}{0.43}
\newcommand{\ratNflameThirty}{3}
\newcommand{\ratBSRfltrustTen}{0.00}  % mean of 3 poison rates
\newcommand{\ratDRfltrustTen}{0.00}
\newcommand{\ratNfltrustTen}{3}
\newcommand{\ratBSRfltrustTwenty}{0.00}  % mean of 3 poison rates
\newcommand{\ratDRfltrustTwenty}{0.00}
\newcommand{\ratNfltrustTwenty}{3}
\newcommand{\ratBSRfltrustThirty}{0.00}  % mean of 3 poison rates
\newcommand{\ratDRfltrustThirty}{0.00}
\newcommand{\ratNfltrustThirty}{3}
\newcommand{\ratBSRfoolsgoldTen}{\PENDING{}}  % mean of 0 poison rates
\newcommand{\ratDRfoolsgoldTen}{\PENDING{}}
\newcommand{\ratNfoolsgoldTen}{\PENDING{}}
\newcommand{\ratBSRfoolsgoldTwenty}{0.03}  % mean of 1 poison rates
\newcommand{\ratDRfoolsgoldTwenty}{0.75}
\newcommand{\ratNfoolsgoldTwenty}{1}
\newcommand{\ratBSRfoolsgoldThirty}{\PENDING{}}  % mean of 0 poison rates
\newcommand{\ratDRfoolsgoldThirty}{\PENDING{}}
\newcommand{\ratNfoolsgoldThirty}{\PENDING{}}
\newcommand{\ratBSRgradcosTen}{0.00}  % mean of 3 poison rates
\newcommand{\ratDRgradcosTen}{0.93}
\newcommand{\ratNgradcosTen}{3}
\newcommand{\ratBSRgradcosTwenty}{0.16}  % mean of 3 poison rates
\newcommand{\ratDRgradcosTwenty}{0.85}
\newcommand{\ratNgradcosTwenty}{3}
\newcommand{\ratBSRgradcosThirty}{0.34}  % mean of 3 poison rates
\newcommand{\ratDRgradcosThirty}{0.63}
\newcommand{\ratNgradcosThirty}{3}

% ---- donor triggers (jointly realizable by construction) ---------------
\newcommand{\dnBSRfedavgFixed}{\PENDING{}}
\newcommand{\dnDRfedavgFixed}{\PENDING{}}
\newcommand{\dnBSRfedavgStatistical}{\PENDING{}}
\newcommand{\dnDRfedavgStatistical}{\PENDING{}}
\newcommand{\dnBSRfedavgTemporal}{\PENDING{}}
\newcommand{\dnDRfedavgTemporal}{\PENDING{}}
\newcommand{\dnBSRkrumFixed}{\PENDING{}}
\newcommand{\dnDRkrumFixed}{\PENDING{}}
\newcommand{\dnBSRkrumStatistical}{\PENDING{}}
\newcommand{\dnDRkrumStatistical}{\PENDING{}}
\newcommand{\dnBSRkrumTemporal}{\PENDING{}}
\newcommand{\dnDRkrumTemporal}{\PENDING{}}
\newcommand{\dnBSRflameFixed}{\PENDING{}}
\newcommand{\dnDRflameFixed}{\PENDING{}}
\newcommand{\dnBSRflameStatistical}{\PENDING{}}
\newcommand{\dnDRflameStatistical}{\PENDING{}}
\newcommand{\dnBSRflameTemporal}{\PENDING{}}
\newcommand{\dnDRflameTemporal}{\PENDING{}}
\newcommand{\dnBSRfltrustFixed}{\PENDING{}}
\newcommand{\dnDRfltrustFixed}{\PENDING{}}
\newcommand{\dnBSRfltrustStatistical}{\PENDING{}}
\newcommand{\dnDRfltrustStatistical}{\PENDING{}}
\newcommand{\dnBSRfltrustTemporal}{\PENDING{}}
\newcommand{\dnDRfltrustTemporal}{\PENDING{}}
\newcommand{\dnBSRgradcosFixed}{\PENDING{}}
\newcommand{\dnDRgradcosFixed}{\PENDING{}}
\newcommand{\dnBSRgradcosStatistical}{\PENDING{}}
\newcommand{\dnDRgradcosStatistical}{\PENDING{}}
\newcommand{\dnBSRgradcosTemporal}{\PENDING{}}
\newcommand{\dnDRgradcosTemporal}{\PENDING{}}

% ---- FLNET2023 (single-seed OAT sweep) ---------------------------------
\newcommand{\flCleanCAfedavg}{0.9996}
\newcommand{\flCleanFonefedavg}{0.9995}
\newcommand{\flBSRfedavgFixed}{0.998}
\newcommand{\flDRfedavgFixed}{\PENDING{}}
\newcommand{\flBSRfedavgStatistical}{0.455}
\newcommand{\flDRfedavgStatistical}{\PENDING{}}
\newcommand{\flBSRfedavgTemporal}{0.567}
\newcommand{\flDRfedavgTemporal}{\PENDING{}}
\newcommand{\flCleanCAkrum}{0.9997}
\newcommand{\flCleanFonekrum}{0.9992}
\newcommand{\flBSRkrumFixed}{0.007}
\newcommand{\flDRkrumFixed}{1.00}
\newcommand{\flBSRkrumStatistical}{0.000}
\newcommand{\flDRkrumStatistical}{1.00}
\newcommand{\flBSRkrumTemporal}{0.007}
\newcommand{\flDRkrumTemporal}{1.00}
\newcommand{\flCleanCAflame}{0.8888}
\newcommand{\flCleanFoneflame}{0.5743}
\newcommand{\flBSRflameFixed}{0.017}
\newcommand{\flDRflameFixed}{1.00}
\newcommand{\flBSRflameStatistical}{0.028}
\newcommand{\flDRflameStatistical}{1.00}
\newcommand{\flBSRflameTemporal}{0.019}
\newcommand{\flDRflameTemporal}{1.00}
\newcommand{\flCleanCAfltrust}{0.9642}
\newcommand{\flCleanFonefltrust}{0.5963}
\newcommand{\flBSRfltrustFixed}{0.016}
\newcommand{\flDRfltrustFixed}{0.00}
\newcommand{\flBSRfltrustStatistical}{0.016}
\newcommand{\flDRfltrustStatistical}{0.00}
\newcommand{\flBSRfltrustTemporal}{0.016}
\newcommand{\flDRfltrustTemporal}{0.00}
\newcommand{\flCleanCAfoolsgold}{\PENDING{}}
\newcommand{\flCleanFonefoolsgold}{\PENDING{}}
\newcommand{\flBSRfoolsgoldFixed}{\PENDING{}}
\newcommand{\flDRfoolsgoldFixed}{\PENDING{}}
\newcommand{\flBSRfoolsgoldStatistical}{\PENDING{}}
\newcommand{\flDRfoolsgoldStatistical}{\PENDING{}}
\newcommand{\flBSRfoolsgoldTemporal}{\PENDING{}}
\newcommand{\flDRfoolsgoldTemporal}{\PENDING{}}
\newcommand{\flCleanCAgradcos}{0.9997}
\newcommand{\flCleanFonegradcos}{0.9996}
\newcommand{\flBSRgradcosFixed}{0.981}
\newcommand{\flDRgradcosFixed}{0.07}
\newcommand{\flBSRgradcosStatistical}{0.001}
\newcommand{\flDRgradcosStatistical}{0.53}
\newcommand{\flBSRgradcosTemporal}{0.007}
\newcommand{\flDRgradcosTemporal}{0.38}

% ---- non-IID (Dirichlet) -----------------------------------------------
\newcommand{\niCleanCAfedavg}{0.777}
\newcommand{\niCleanFonefedavg}{0.527}
\newcommand{\niBSRfedavgFixed}{0.792}
\newcommand{\niCAfedavgFixed}{0.765}
\newcommand{\niDRfedavgFixed}{\PENDING{}}
\newcommand{\niBSRfedavgStatistical}{\PENDING{}}
\newcommand{\niCAfedavgStatistical}{\PENDING{}}
\newcommand{\niDRfedavgStatistical}{\PENDING{}}
\newcommand{\niBSRfedavgTemporal}{\PENDING{}}
\newcommand{\niCAfedavgTemporal}{\PENDING{}}
\newcommand{\niDRfedavgTemporal}{\PENDING{}}
\newcommand{\niCleanCAkrum}{0.670}
\newcommand{\niCleanFonekrum}{0.524}
\newcommand{\niBSRkrumFixed}{0.926}
\newcommand{\niCAkrumFixed}{0.636}
\newcommand{\niDRkrumFixed}{0.00}
\newcommand{\niBSRkrumStatistical}{\PENDING{}}
\newcommand{\niCAkrumStatistical}{\PENDING{}}
\newcommand{\niDRkrumStatistical}{\PENDING{}}
\newcommand{\niBSRkrumTemporal}{\PENDING{}}
\newcommand{\niCAkrumTemporal}{\PENDING{}}
\newcommand{\niDRkrumTemporal}{\PENDING{}}
\newcommand{\niCleanCAflame}{0.519}
\newcommand{\niCleanFoneflame}{0.292}
\newcommand{\niBSRflameFixed}{0.014}
\newcommand{\niCAflameFixed}{0.559}
\newcommand{\niDRflameFixed}{0.57}
\newcommand{\niBSRflameStatistical}{\PENDING{}}
\newcommand{\niCAflameStatistical}{\PENDING{}}
\newcommand{\niDRflameStatistical}{\PENDING{}}
\newcommand{\niBSRflameTemporal}{\PENDING{}}
\newcommand{\niCAflameTemporal}{\PENDING{}}
\newcommand{\niDRflameTemporal}{\PENDING{}}
\newcommand{\niCleanCAfltrust}{0.508}
\newcommand{\niCleanFonefltrust}{0.415}
\newcommand{\niBSRfltrustFixed}{0.070}
\newcommand{\niCAfltrustFixed}{0.503}
\newcommand{\niDRfltrustFixed}{0.00}
\newcommand{\niBSRfltrustStatistical}{\PENDING{}}
\newcommand{\niCAfltrustStatistical}{\PENDING{}}
\newcommand{\niDRfltrustStatistical}{\PENDING{}}
\newcommand{\niBSRfltrustTemporal}{\PENDING{}}
\newcommand{\niCAfltrustTemporal}{\PENDING{}}
\newcommand{\niDRfltrustTemporal}{\PENDING{}}
\newcommand{\niCleanCAfoolsgold}{\PENDING{}}
\newcommand{\niCleanFonefoolsgold}{\PENDING{}}
\newcommand{\niBSRfoolsgoldFixed}{\PENDING{}}
\newcommand{\niCAfoolsgoldFixed}{\PENDING{}}
\newcommand{\niDRfoolsgoldFixed}{\PENDING{}}
\newcommand{\niBSRfoolsgoldStatistical}{\PENDING{}}
\newcommand{\niCAfoolsgoldStatistical}{\PENDING{}}
\newcommand{\niDRfoolsgoldStatistical}{\PENDING{}}
\newcommand{\niBSRfoolsgoldTemporal}{\PENDING{}}
\newcommand{\niCAfoolsgoldTemporal}{\PENDING{}}
\newcommand{\niDRfoolsgoldTemporal}{\PENDING{}}
\newcommand{\niCleanCAgradcos}{0.784}
\newcommand{\niCleanFonegradcos}{0.531}
\newcommand{\niBSRgradcosFixed}{0.004}
\newcommand{\niCAgradcosFixed}{0.773}
\newcommand{\niDRgradcosFixed}{0.47}
\newcommand{\niBSRgradcosStatistical}{\PENDING{}}
\newcommand{\niCAgradcosStatistical}{\PENDING{}}
\newcommand{\niDRgradcosStatistical}{\PENDING{}}
\newcommand{\niBSRgradcosTemporal}{\PENDING{}}
\newcommand{\niCAgradcosTemporal}{\PENDING{}}
\newcommand{\niDRgradcosTemporal}{\PENDING{}}

% ---- cell counts (provenance) ------------------------------------------
\newcommand{\numCellsSingle}{200}  % single-seed cells in summary.csv
\newcommand{\numCellsMulti}{28}  % grouped cells in summary_multiseed.csv
\newcommand{\numSeeds}{5}


\journal{Journal of Information Security and Applications}

\begin{document}

\begin{frontmatter}

\title{Direction Is Not Enough: Trigger Salience and Adaptive Collusion
against Byzantine-Robust Federated Intrusion Detection}

%---------------------------------------------------------------------
% AUTHOR BLOCK
%---------------------------------------------------------------------
\author[mcs]{Usama Mustafa\corref{cor1}}
\ead{usamamustafa1@gmail.com}

\author[mcs]{Imran Rashid}
\ead{irashid@mcs.edu.pk}

\cortext[cor1]{Corresponding author.}

\affiliation[mcs]{organization={Department of Information Security,
    Military College of Signals, National University of Sciences and
    Technology},
    city={Rawalpindi},
    country={Pakistan}}

%---------------------------------------------------------------------
\begin{abstract}
Federated learning trains a shared intrusion detector without pooling raw traffic, opening
the aggregation step to poisoning by compromised participants. We study backdoor attacks on
a federated CNN\,+\,LSTM network-flow classifier and the Byzantine-robust defenses meant to
stop them, on two 2023 benchmarks (CICIoT2023, FLNET2023) under IID and non-IID data. Our findings are mostly negative and concern the \emph{signal} these defenses read, not
any one defense. First, whether a similarity-based rule holds is decided by how far the trigger
rotates the poisoned update off the honest consensus direction, measured per round rather
than inferred from the trigger's name, and not by the malicious-client fraction. Second,
that signal is cheap to remove: zero-sum noise across the colluding clients decorrelates
them while leaving the aggregate untouched, collapsing their mean pairwise similarity from
% FULL-SCALE CHECK: the "below 0.35" figure below comes from the reduced-scale
% validation run (see the provenance note in Sec. 7.4). tools/check_manuscript.py
% compares it against the adMalCos* macros as lane 1a lands; at the first
% full-scale cell the worst measured value is 0.428, i.e. ABOVE the claim.
% Revisit this sentence -- and the matching one in the introduction -- when the
% lane completes. Do not leave it as a bare literal.
$0.99$ to below $0.35$ at no cost to the backdoor. Every direction-based rule here is
evaluated against it. Third, detection rates reported without false-positive rates mislead:
FLAME reaches a backdoor rate of $\msBSRflameFixed$ while flagging $\msFPRpctflameFixed\%$ of \emph{honest} clients,
which is where its $\msCAdeficitFlame$ clean-accuracy deficit comes from. Fourth, a joint-distribution audit shows this
literature's ``in-distribution'' premise does not survive the data: a trigger deviating by
$0.5\sigma$ on any single feature can still be further from the benign mean than all
but $\rzExceedFixedQtr$ of $\rzRefRows$ real benign flows. We propose \textsc{GradCos}, a
reference-free aggregator that ties the strongest baseline on conspicuous triggers with no
clean server data and at $+\msCAadvGradcos$ clean accuracy, and report its blind spot, seed bimodality
and failure against that attacker.
\end{abstract}

\begin{keyword}
federated learning \sep backdoor attacks \sep intrusion detection \sep
Byzantine-robust aggregation \sep adaptive attacks \sep IoT security \sep
poisoning
\end{keyword}

\end{frontmatter}

\section{Introduction}

Intrusion detection on IoT and edge networks is a natural fit for federated
learning (FL). The traffic a detector needs is exactly the traffic operators are
least willing to centralize, and FL~\cite{mcmahan2017,kairouz2021} keeps it on
the device: each participant trains locally and shares only model updates, which
a server averages into a global detector. The catch is that the averaging step
trusts whoever shows up. A compromised participant---a common assumption in a
large, heterogeneous IoT fleet---can submit crafted updates that steer the shared
model. Among poisoning attacks, \emph{backdoors} are the most insidious: rather
than degrade the model outright, they teach it a hidden rule (``traffic carrying
pattern $\tau$ is benign'') while leaving normal behavior
intact~\cite{bagdasaryan2020,xie2020}. In an intrusion detector this lets an
attacker wave their own traffic past the model, and nothing in the aggregate
accuracy gives it away.

A sizeable literature answers this with Byzantine-robust aggregation. Krum picks
the update in the densest neighbourhood~\cite{blanchard2017}; FLTrust scores each
update against a direction derived from a small clean dataset~\cite{cao2021fltrust};
FLAME clusters updates, clips their norms, and adds noise~\cite{nguyen2022flame};
FoolsGold penalizes clients whose update histories are too similar to each
other~\cite{fung2020foolsgold}. Strip away the differences and almost all of them
read the same underlying quantity: \emph{where an update points}. That is the
observation this paper is organized around, and most of what we report follows
from taking it seriously.

We build a federated CNN\,+\,LSTM flow classifier, attack it with three families
of tabular flow-feature triggers, and defend it with ten aggregation rules---five
that read update direction, five that read only magnitude or coordinate order
statistics. We instrument the server to record, every round, the geometry each
defense is actually looking at. Working through CICIoT2023 and FLNET2023, IID and
non-IID, with the headline comparison over five seeds, we reach a picture that is
more negative and, we think, more useful than a ranking of defenses.

\begin{itemize}
\item \textbf{Salience, measured rather than named.} We define the salience of a
trigger operationally, as the gap between how well an honest update and a
poisoned update align with the honest consensus direction, and log it per round.
Robustness tracks that quantity, not the attacker count: conspicuous triggers are
filtered by every direction-based defense even at $30\%$ malicious clients, while
triggers that barely rotate the update defeat all of them. Sweeping trigger
magnitude at a fixed feature set turns this from a contrast between two groups of
triggers into a dose--response along one measured axis.

\item \textbf{The signal is cheap to remove.} Direction-based defenses detect
colluders because colluders resemble each other. An attacker who adds independent
per-client noise, mean-centred across the malicious cohort so it cancels under
averaging, keeps the backdoor and destroys the resemblance---mean pairwise cosine
% FULL-SCALE CHECK: the "below 0.35" figure below comes from the reduced-scale
% validation run (see the provenance note in Sec. 7.4). tools/check_manuscript.py
% compares it against the adMalCos* macros as lane 1a lands; at the first
% full-scale cell the worst measured value is 0.428, i.e. ABOVE the claim.
% Revisit this sentence -- and the matching one in the introduction -- when the
% lane completes. Do not leave it as a bare literal.
among colluders falls from $0.99$ to below $0.35$ while the honest cohort is unchanged
(measured directly; Sec.~\ref{sec:adaptive}). Norm-matching the result hides the
residual magnitude tell. We evaluate this attack against every direction-based
defense including our own.

\item \textbf{GradCos}, a reference-free aggregator that identifies the
colluding minority from the pairwise cosine structure of the updates, needing no
clean server dataset (unlike FLTrust) and no estimate of the attacker count
(unlike Krum). On conspicuous triggers it ties FLTrust's backdoor rate ($\msBSRfltrustFixed$)
while keeping $+\msCAadvGradcos$ clean accuracy and holding where Krum collapses. We are
explicit about the rest: it fails on low-salience triggers exactly as its peers
do, it is bimodal across seeds on our second dataset, and the adaptive attack
above defeats it. Its contribution is the diagnosis it enables, not a ranking.

\item \textbf{Detection rates without false-positive rates are not evidence.}
FLAME reaches a backdoor rate of $\msBSRflameFixed$ on the fixed trigger while flagging
$0.40\pm0.04$ of honest clients; that is where its clean accuracy goes. On
low-salience triggers \textsc{GradCos} flags honest clients marginally more often
than malicious ones. Neither fact is visible in a backdoor-success-only
comparison, and both change what the numbers mean.

\item \textbf{A realizability audit of the triggers themselves.} Marginal
$\sigma$ checks, standard in this literature, say almost nothing about whether a
stamped flow could exist. We measure joint plausibility directly and find our own
threat model wanting in a specific, reportable way, and we supply a
donor-flow trigger construction that fixes part of it.
\end{itemize}

We release the code, configurations, per-round update-geometry logs, and the
audit so the grid can be regenerated with a single command.

\section{Related Work}

\textbf{Backdoors in federated learning.} Bagdasaryan et
al.~\cite{bagdasaryan2020} showed that model replacement lets a single client
implant a durable backdoor; Xie et al.~\cite{xie2020} distributed the trigger
across colluding clients (DBA) so each update looks more benign. Wang et
al.~\cite{wang2020tails} argued backdoors on rare sub-populations are especially
hard to remove, while Sun et al.~\cite{sun2019} found norm clipping and noise
blunt many backdoors. Most of this work uses image classification with patch
triggers. Intrusion detection differs in that the feature space is tabular flow
statistics whose coordinates are jointly constrained, which is a problem for the
trigger design rather than a detail of it (Sec.~\ref{sec:realizable}).

\textbf{Byzantine-robust aggregation.} Krum and Multi-Krum~\cite{blanchard2017}
give distance-based robustness; coordinate-wise median and trimmed
mean~\cite{yin2018}, Bulyan~\cite{mhamdi2018bulyan} and the geometric median
(RFA)~\cite{pillutla2022rfa} give order-statistic robustness.
FLTrust~\cite{cao2021fltrust} adds a trusted root dataset to anchor a reference
direction. FoolsGold~\cite{fung2020foolsgold} targets sybil collusion by
penalizing clients whose updates are too similar to each other---the opposite
reading of the same cosine signal we use. FLAME~\cite{nguyen2022flame} and
DeepSight~\cite{rieger2022deepsight} combine clustering with clipping/filtering.
We evaluate both families side by side, because the separation between them is
what our main finding turns on.

\textbf{Adaptive attacks.} Shejwalkar and Houmansadr~\cite{shejwalkar2021} and
Fang et al.~\cite{fang2020local} showed that defenses evaluated only against
static attackers routinely fall to attackers optimized against them; Baruch et
al.~\cite{baruch2019little} showed how little perturbation is needed when the
attacker stays inside the honest variance. The closest work to ours is
3DFed~\cite{li20233dfed}, which evades DeepSight, FoolsGold, FLAME and
FLDetector~\cite{zhang2022fldetector} by adapting its parameters from feedback
about the previous global model, and the learning-based framework of Yang et
al.~\cite{yang2024learning}, which trains the attack against the defense.
Our evasion (Sec.~\ref{sec:adaptive}) is a deliberately unsophisticated instance
of the same lesson---no optimization, no feedback channel, one scalar
knob---which is what makes it awkward for the defenses it breaks: a weaker
attacker that still works is worse news for a defense than a stronger one that
does. Read alongside Fang et al.~\cite{fang2025foundation}, who argue the field
should harden established aggregation rules rather than keep proposing new ones,
our results say something about \emph{which} established rules are worth
hardening: not the ones reading update direction.
Pierazzi et al.~\cite{pierazzi2020problem} make the problem-space point that
motivates our realizability audit: a feature-space perturbation is only an attack
if some real object realizes it.

\textbf{Federated IDS.} Federated detectors have been built on modern IoT
benchmarks such as CICIoT2023~\cite{neto2023ciciot} and the federated traffic
benchmark FLNET2023~\cite{flnet2023}. Poisoning of these detectors specifically
has been studied by Zhang et al.~\cite{zhang2022secfednids}, who defend at both
the model and the data level; by Merzouk et al.~\cite{merzouk2023param}, who
parameterize the attacker's budget; and by Sameera et
al.~\cite{sameera2025weidetect}, who fit validation scores to a Weibull
distribution to exclude poisoned models. All three evaluate against static
attackers, and none reports the false-positive cost of the exclusions its
detector makes. Comparatively little of this work stress-tests the aggregation
layer against stealthy backdoors with clean-accuracy and false-positive
accounting, which is the gap this paper addresses.

\section{Threat Model and System Setup}

\textbf{Federated detector.} $N=20$ clients hold private shards of a
network-flow dataset and train a shared classifier for $R=100$ synchronous
rounds at full participation, one local epoch per round (batch size $256$, Adam,
learning rate $10^{-3}$). The model treats the $F$-dimensional standardized
feature vector as a length-$F$, single-channel sequence: a Conv1D layer ($64$
filters, kernel $3$, ReLU, BatchNorm) extracts local feature interactions, an
LSTM ($128$ hidden units) reads the resulting sequence, and a dense head ($64$
units, dropout $0.3$) with softmax produces the class posterior. Aggregation is
performed by genuine Flower~\cite{beutel2020flower} \texttt{Strategy} objects
driven by an in-process round loop, so every defense receives real per-client
\texttt{FitRes} updates.

\textbf{Adversary.} A fraction $m\in\{0.1,0.2,0.3\}$ of clients is malicious and
colludes on a single backdoor objective: make trigger-carrying attack traffic be
classified as benign. Each malicious client poisons a fraction
$p\in\{0.05,0.1,0.2\}$ of its \emph{attack} samples by stamping the trigger and
relabelling them benign, then trains normally. The adversary knows the model
architecture and the trigger but controls neither the server nor honest clients.

We are explicit that this is a \emph{feature-space} threat model: the trigger is
a perturbation of the standardized feature vector the classifier consumes, and we
do not claim a corresponding packet-level realization.
Sec.~\ref{sec:realizable} measures how far that abstraction is from realizable
traffic, because we found the usual justification for it does not hold.

\textbf{Adaptive adversary.} Beyond static poisoning we consider an attacker who
knows the defense. Sec.~\ref{sec:adaptive} gives an evasion that targets the one
property every direction-based rule here depends on---that colluders resemble one
another---at the cost of one scalar parameter and no optimization.

\textbf{Server capability.} The server sees only model updates and has no access
to raw client data. FLTrust is the one deliberate exception: it is granted a
small class-balanced clean root set ($\approx$300 samples), so its results should
be read as an upper bound assuming a resource the other defenses do without.

\textbf{Metrics.} We report clean accuracy (CA) and macro-F1 on non-triggered
traffic; backdoor success rate (BSR), the fraction of triggered attack samples
called benign; and, for defenses that flag clients, detection rate (DR) and
false-positive rate (FPR) against the known malicious set. \textbf{DR and FPR are
always reported together.} A defense that flags clients converts a false positive
directly into lost training signal, so a detection rate quoted alone says nothing
about what the detector cost; Sec.~\ref{sec:fpr} shows this is not a hypothetical
concern for two of the baselines. Defenses that never single out a client
(FedAvg and the coordinate-wise rules) report DR/FPR as N/A rather than zero. We
also log the per-round update geometry defined next.

\textbf{Salience, defined.} Let $\delta_i$ be client $i$'s update relative to the
round's global model and $\bar\delta_h$ the mean honest update. We define the
\emph{salience gap}
\begin{equation}
\label{eq:salience}
S \;=\; \frac{1}{|H|}\sum_{i\in H}\cos(\delta_i,\bar\delta_h)
      \;-\; \frac{1}{|M|}\sum_{j\in M}\cos(\delta_j,\bar\delta_h),
\end{equation}
with $H$ and $M$ the honest and malicious sets. $S$ is $0$ when the backdoor does
not rotate the update at all and grows as the trigger becomes more conspicuous.
It is exactly the quantity a cosine filter must resolve, it is measured rather
than assumed, and it lets us regress outcomes on a continuous axis instead of on
a trigger's name. We log $S$ every round, together with the mean off-diagonal
cosine within each group (the collusion signal) and the ratio of median malicious
to median honest update norm (what a magnitude rule sees).

\section{Trigger Design and Realizability}
\label{sec:realizable}

Because the classifier consumes standardized features, we define triggers in
normalized space, so train-time poisoning and test-time evaluation are consistent
by construction. Table~\ref{tab:trig} lists the three families: flag/header
features, volume/dispersion features, and inter-arrival/count features. Each is
mixed-sign at roughly $1.5$--$2\sigma$ and verified to have a clean-model BSR near
zero, so the association must be learned from poisoned updates rather than being
something the unpoisoned model already does. Every clean cell in our grid now
reports that clean-model BSR for the whole trigger battery, so the check is a
per-run control rather than a one-off.

\begin{table}[t]
\caption{Trigger families (values in units of $\sigma$; CICIoT2023 schema).}
\label{tab:trig}
\centering
\footnotesize
\setlength{\tabcolsep}{4pt}
\begin{tabular}{@{}lll@{}}
\toprule
Trigger & Semantics & Example features / values \\
\midrule
Fixed & header \& TCP flags & TTL$=+2$, hdr\_len$=-1.5$, syn$=+2$ \\
Statistical & volume / dispersion & Rate$=-1.5$, Std$=+2$, Tot\,size$=+1.5$ \\
Temporal & inter-arrival / counts & IAT$=+2$, Number$=+1.5$, size$=-1.5$ \\
\bottomrule
\end{tabular}
\end{table}

\subsection{Marginal checks do not establish realizability}

The clean-model BSR control and the $1.5$--$2\sigma$ magnitudes are the evidence
usually offered for calling such a trigger ``in-distribution'' and therefore
stealthy. Both are marginal, and flow features are not independent: rate,
duration, total size, packet count and inter-arrival time are functions of one
another, so a vector whose every coordinate is individually ordinary can be a
combination the network could never emit. We therefore measure joint
plausibility directly, on three axes: the Mahalanobis distance of the stamped
sub-vector under the benign covariance restricted to the trigger columns, read
against the same statistic on real benign flows; the distance from the stamped
block to its nearest real flow; and the same nearest-real-flow distance for
fully stamped attack vectors over all $39$ features. Table~\ref{tab:realize}
reports the result on CICIoT2023; the corresponding FLNET2023 audit is
% RETIRES WITH: FLNET realizability audit (last step of run_remediation.sh)
outstanding, so the conclusions below are stated for CICIoT2023 only. \PENDING

\begin{table}[t]
\caption{\textbf{Joint plausibility of the triggers} (CICIoT2023, full training
split). $|\sigma|_{\max}$ is the largest marginal deviation---the number usually
reported. $d_M$ is the joint Mahalanobis distance of the stamped block under the
benign covariance; benign flows themselves have $d_M^{p99}$ as shown. $d_{NN}$ is
the distance from a fully stamped attack vector to the nearest real flow, against
$\rzUnstamped$ for a real attack flow's nearest real neighbour. Marginal and joint
plausibility are close to unrelated.}
\label{tab:realize}
\centering
\footnotesize
\begin{tabular}{@{}llrrrr@{}}
\toprule
Trigger & Mode & $|\sigma|_{\max}$ & $d_M$ & $d_M^{p99}$ & $d_{NN}$ \\
\midrule
Fixed       & $\sigma$, $0.25\times$ & \rzSigFixedQtr & $\rzDMFixedQtr$ & \rzPnnFixed & \rzNNFixedQtr \\
Fixed       & $\sigma$, $1\times$    & \rzSigFixedOne & $\rzDMFixedOne$ & \rzPnnFixed & \rzNNFixedOne \\
Fixed       & $\sigma$, $3\times$    & \rzSigFixedThree & $\rzDMFixedThree$ & \rzPnnFixed & \rzNNFixedThree \\
Fixed       & donor                  & \rzSigFixedDonor & $\textbf{\rzDMFixedDonor}$ & \rzPnnFixed & \rzNNFixedDonor \\
\midrule
Statistical & $\sigma$, $0.25\times$ & \rzSigStatisticalQtr & $\textbf{\rzDMStatisticalQtr}$ & \rzPnnStatistical & \rzNNStatisticalQtr \\
Statistical & $\sigma$, $1\times$    & \rzSigStatisticalOne & $\textbf{\rzDMStatisticalOne}$ & \rzPnnStatistical & \rzNNStatisticalOne \\
Statistical & $\sigma$, $3\times$    & \rzSigStatisticalThree & $\textbf{\rzDMStatisticalThree}$ & \rzPnnStatistical & \rzNNStatisticalThree \\
Statistical & donor                  & \rzSigStatisticalDonor & $\textbf{\rzDMStatisticalDonor}$ & \rzPnnStatistical & \rzNNStatisticalDonor \\
\midrule
Temporal    & $\sigma$, $0.25\times$ & \rzSigTemporalQtr & $\rzDMTemporalQtr$ & \rzPnnTemporal & \rzNNTemporalQtr \\
Temporal    & $\sigma$, $1\times$    & \rzSigTemporalOne & $\rzDMTemporalOne$ & \rzPnnTemporal & \rzNNTemporalOne \\
Temporal    & $\sigma$, $3\times$    & \rzSigTemporalThree & $\rzDMTemporalThree$ & \rzPnnTemporal & \rzNNTemporalThree \\
Temporal    & donor                  & \rzSigTemporalDonor & $\textbf{\rzDMTemporalDonor}$ & \rzPnnTemporal & \rzNNTemporalDonor \\
\bottomrule
\end{tabular}
\end{table}

Three things follow, and the first two are corrections to our own premise.

\emph{Marginal magnitude is a poor proxy for joint plausibility.} The fixed
trigger at $0.25\times$ has a largest marginal deviation of
$0.5\sigma$---by the usual standard, unremarkable---yet sits at
$d_M=\rzDMFixedQtr$ against a benign $99$th percentile of $\rzPnnFixed$: of the
$\rzRefRows$ real benign flows the audit uses as its reference, only
$\rzExceedFixedQtr$ sit further from the benign mean than this trigger does. We quote
the count rather than a percentile because at that tail the percentile is carried by a
handful of rows and reads more precisely than the sample supports. Conversely the donor triggers have marginals of
$\rzSigTemporalDonor$--$\rzSigStatisticalDonor\sigma$ and are jointly
\emph{typical} ($d_M=\rzDMTemporalDonor$--$\rzDMStatisticalDonor$), because they are copied from
real flows. The two quantities are close to independent, so reporting the first
and claiming the second, as we and others have done, is not sound.

\emph{The three triggers are not jointly comparable.} The statistical trigger is
genuinely inside the benign joint envelope at every scale we test
($d_M=\rzDMStatisticalQtr$--$\rzDMStatisticalThree$ against
$d_M^{p99}=\rzPnnStatistical$). The temporal trigger is not, by
more than two orders of magnitude ($d_M\approx10^3$ against $\rzPnnTemporal$): setting IAT,
packet count and total size independently produces a flow whose timing and volume
contradict each other. The paper's grouping of ``statistical and temporal'' as a
single low-salience family is therefore sound with respect to update geometry but
\emph{not} with respect to realizability, and results resting on the temporal
trigger should be read as feature-space results only. We flag this wherever the
temporal trigger drives a conclusion.

\emph{A partial fix.} Our \texttt{donor} construction copies the trigger block
from a single real benign flow, chosen to be distinctive relative to attack
traffic while remaining inside the benign joint envelope (below its $95$th
Mahalanobis percentile). By construction the stamped block is a combination the
network actually produced, and it also matches a more realistic adversary---one
who shapes traffic to imitate a benign flow rather than dialing individual
statistics to chosen $\sigma$ values. It brings all three triggers to
$d_M\approx\rzDMTemporalDonor$--$\rzDMStatisticalDonor$. It does not close the
gap: $d_{NN}$ stays at $\rzNNFixedDonor$--$\rzNNStatisticalDonor$, because
pasting a real benign block into a real attack flow still yields a cross-
combination that no single record exhibits. Full problem-space
realizability~\cite{pierazzi2020problem} would require generating traffic and
re-extracting features, which is outside this paper's scope. We report the
donor-trigger results alongside the $\sigma$-trigger results and treat the
difference as part of the finding.

\textbf{Realizability and backdoor-ness pull against each other.} Constructing
the donor trigger exposed a constraint we had not anticipated, and it is the more
interesting half of this section. Selecting the \emph{most} distinctive benign
donor---the row least like attack traffic on the trigger columns---is close to
selecting the \emph{maximally benign-looking} block, and a flow stamped with it
is called benign by the \emph{clean} model with no poisoning at all. Our
clean-model BSR control catches this immediately, which is what the control is
for. But it means such a trigger is an evasion rather than a backdoor, and BSR
stops measuring a learned association. Pushing the other way---choosing a donor
the clean model does not already accept---makes the pattern less separable from
attack traffic and correspondingly harder for the model to key on. We take the
median of the distinctiveness ordering rather than its extreme, and screen every
donor trigger by its clean-model BSR before reading any result from it. That
screen passes: against an unpoisoned model the median-quantile donors give
clean-model BSR $0.006$ (fixed), $0.000$ (statistical) and $0.008$ (temporal),
against $0.006$/$0.000$/$0.001$ for the corresponding $\sigma$ triggers---so the
donor triggers are backdoor candidates on the same footing, where the
maximum-distinctiveness donor for statistical had sat at $1.000$.

We do not think this tension is an artifact of our construction. Any trigger that
is simultaneously realizable as benign-looking traffic and learnable as a rule
distinct from ``benign'' is being asked to be two things at once, and the
narrowness of the window between them is a constraint on problem-space backdoors
in network IDS that the feature-space formulation hides entirely. We flag it as
such, and note that it is only visible because the clean-model BSR control runs
per cell rather than once.

\section{Defenses}

We evaluate ten aggregation rules in two families. All subclass FedAvg and
differ only in \texttt{aggregate\_fit}. The split between the families is the
paper's organising distinction, not a presentational one.

\subsection{Direction-based rules}

These decide what to keep from where an update \emph{points}.

\textbf{FedAvg}~\cite{mcmahan2017} is the undefended, example-weighted mean---the
attack's baseline.

\textbf{Krum / Multi-Krum}~\cite{blanchard2017} scores each update by the sum of
squared distances to its $n-f-2$ nearest neighbours and keeps the update(s) in
the densest region; we use Multi-Krum with $f$ clamped to what the cohort size
supports.

\textbf{FLTrust}~\cite{cao2021fltrust} trains the current global model one step
on the server's clean root set to get a trusted direction $g_0$, weights each
client update by $\mathrm{ReLU}(\cos(g_i,g_0))$, rescales it to $\|g_0\|$, and
averages.

\textbf{FLAME}~\cite{nguyen2022flame} clusters updates by cosine distance
(HDBSCAN), keeps the majority cluster, clips admitted updates to their median
norm, and adds Gaussian noise. We fixed a bug that had invalidated our early
FLAME runs: the differential-privacy noise was reaching BatchNorm
\texttt{running\_var} buffers and driving them negative, so eval-mode BatchNorm
computed $\sqrt{\cdot}$ of a negative number and produced NaN predictions. The
noise is now masked to trainable parameters only. We also now seed FLAME's noise
from the run seed rather than a constant, so its multi-seed dispersion is real
rather than an artifact of a fixed noise draw.

\textbf{FoolsGold}~\cite{fung2020foolsgold} reads the same cosine signal in
reverse: it accumulates each client's historical update direction and
down-weights clients whose histories are too similar to one another.

\subsection{Magnitude and order-statistic rules}

These never inspect direction, which makes them the control group for our main
claim. If low-salience triggers evade defenses because they do not move the
update direction, then rules that do not read direction should not share the
blind spot for that reason.

\textbf{Coordinate-wise median} and \textbf{trimmed mean}~\cite{yin2018} take
per-coordinate order statistics. \textbf{RFA}~\cite{pillutla2022rfa} computes the
geometric median via smoothed Weiszfeld iterations---a genuine multivariate
median, so unlike the coordinate-wise rules it is not defeated by an attacker who
spreads a perturbation across coordinates. \textbf{Bulyan}~\cite{mhamdi2018bulyan}
combines iterated Multi-Krum selection with a coordinate-wise trimmed mean.
\textbf{Norm-clip}~\cite{sun2019} clips each update to the median norm and adds
Gaussian noise; it is FLAME with the cosine stage removed, which isolates how
much of FLAME's behavior is the magnitude stage.

\subsection{GradCos (proposed)}

The premise is that honest clients optimize the same objective on
similar-distribution data, so their update directions are mutually well-aligned,
whereas a backdoor pulls a client's update toward the trigger\,$\rightarrow$\,benign
association. Crucially the signal is purely inter-client: no clean server dataset
is needed. Let $\delta_i = w_i - w^{(t)}$ and $\hat\delta_i =
\delta_i/\|\delta_i\|$; the method works from $C_{ij} = \langle
\hat\delta_i,\hat\delta_j\rangle$.

The naive reading---flag clients whose \emph{mean} similarity to the cohort is a
robust-$z$ below the median (our \texttt{meansim} variant)---fails under
collusion: when $20$--$30\%$ of clients submit mutually similar backdoor updates,
each attacker's mean similarity is propped up by its fellow attackers. Our
default \texttt{clique} variant turns collusion into a signal instead. Near
convergence honest updates shrink toward noise and point in diverse directions,
so each honest client's \emph{maximum} similarity to any peer is low; colluding
attackers keep injecting the same direction, so each has a near-twin. GradCos
computes $s_i=\max_{j\neq i} C_{ij}$, flags clients with $s_i > \mathrm{med}(s) +
\kappa\,1.4826\,\mathrm{MAD}(s)$, and averages the rest. A homogeneity gate keeps
everyone when most clients have a twin (early training, or a benign IID cohort
where all updates align); when no clique is found the rule falls back to the
mean-similarity outlier test; a benign-majority cap and safety floor retain at
least half the cohort.

We note plainly what this does and does not contribute. The cosine signal is not
new---FoolsGold reads the same quantity with the opposite sign---and the clique
rule is a different decision procedure on it, not a different signal. Its
practical claim is narrow: on conspicuous triggers it matches the best baseline
without needing clean server data or an attacker-count estimate, and it does so
without the clean-accuracy loss those baselines incur. Sec.~\ref{sec:adaptive}
shows the same clique structure is what makes it evadable.

\textbf{Residual variant.} The full-vector cosine is dominated by the shared
main-task gradient, which is why low-salience triggers slip past it. A
\texttt{residual} mode mean-centres the updates first: $r_i=\delta_i-\bar\delta$,
then runs the same clique rule on $\hat r_i$. Subtracting the cohort mean removes
the common task direction and leaves the colluders' shared component exposed
(Sec.~\ref{sec:residual}).

\subsection{Hyperparameter selection}
\label{sec:tuning}

Several comparisons in this paper turn on hyperparameter-sensitive numbers.
FLAME's clean accuracy depends strongly on its DP noise scale; Krum's and
Bulyan's on the assumed attacker count $f$; the trimmed mean's on $\beta$;
GradCos's on $\kappa$. If those were chosen by inspecting the evaluation grid,
then both our ``no free lunch'' claim about FLAME and GradCos's own result would
be partly artifacts of tuning---and GradCos carries several hand-designed guards
that would be easy to fit to the test conditions. We therefore fix all of them by
one declared protocol, run before the evaluation grid
(\texttt{scripts/calibrate\_defenses.py}):

\begin{itemize}
\item a calibration data split drawn with its own subsample seed, so the rows are
not the evaluation rows;
\item calibration seed $7$, outside the $42$--$46$ evaluation band, so no
calibration draw is also a reported draw;
\item a reduced budget ($25$ rounds, $80$k rows), since the task is ranking
settings against one another, not measuring a final number;
\item one objective declared in advance---maximize $\mathrm{CA}-\mathrm{BSR}$
under attack at a single operating point---which rewards a setting only when
suppressing the backdoor does not cost the model;
\item a single trigger (fixed) for calibration, deliberately: tuning per trigger
would be tuning on the axis under study.
\end{itemize}

Every defense including our own goes through the identical procedure, and the
full grid behind each choice---not just the winner---is released in
\texttt{results/calibration.json}, so the sensitivity of each result to its
% RETIRES WITH: the calibration run (scripts/calibrate_defenses.py)
hyperparameter is visible rather than asserted. \PENDING{} the calibration run
itself; the numbers below use the pre-registered defaults recorded in
\texttt{configs/default.yaml}, and will be regenerated against the calibrated
values.

\section{Experimental Setup}

CICIoT2023~\cite{neto2023ciciot} is used in its 8-class grouped form ($39$
features), stratified-subsampled to $400$k training and $100$k test rows to keep
the $20$-client, $100$-round grid tractable; FLNET2023~\cite{flnet2023} uses the
CICFlowMeter feature schema with category labels. Features are standardized with
a scaler fit on the training split. Clients are partitioned IID, or non-IID via a
Dirichlet split. The CICIoT2023 IID breadth grid is $5$ defenses $\times$
$28$ configurations $=140$ single-seed cells; FLNET2023 is a one-at-a-time sweep
($40$ cells) and the non-IID slice is one configuration across five defenses
($10$ cells). The \textbf{headline operating point} ($m{=}0.2$, $p{=}0.1$) is run
over \textbf{five seeds} (42--46) for all three triggers and six defenses and
reported as mean\,$\pm$\,s.d. Every number quoted in this paper---in the prose and in
every table---is a macro generated from the committed result JSONs by
\texttt{paper/make\_numbers.py}; nothing is hand-entered, and the generator
refuses to emit a number for a cell that has not been run.

We are explicit about which claims rest on which evidence, because the two are
not interchangeable and our own variance results (Sec.~\ref{sec:seed}) show why.
\textbf{Multi-seed evidence} supports the headline defense comparison, the
DR/FPR analysis, the salience finding, and the adaptive-attack result.
\textbf{Single-seed evidence} supports the breadth grid, the malicious-ratio
sweep, the FLNET slice, and the non-IID slice; those are labeled single-seed in
every caption, and we draw only directional conclusions from them.

\section{Results}

\subsection{The attack is invisible to accuracy monitoring}

Under undefended FedAvg the backdoor succeeds everywhere. At the headline
operating point, averaged over five seeds, FedAvg's BSR is
$\msBSRfedavgFixed\pm\msBSRsdfedavgFixed$ on the fixed trigger and
$\msBSRfedavgStatistical$ on both low-salience triggers (Table~\ref{tab:main}),
while its clean accuracy ($\msCleanCAfedavg$) is indistinguishable from the
no-attack model.
Across the broader single-seed grid every one of the $\gridAttackCells$ attack
cells has BSR $\geq\gridMinBSRfedavg$. The poisoning leaves no accuracy footprint, which is what makes it
dangerous: an operator watching validation accuracy sees a healthy detector. This
motivates aggregation-level defenses rather than accuracy monitoring.

\subsection{Headline comparison, with false positives}
\label{sec:fpr}

Table~\ref{tab:main} is the paper's central table, and it now carries FPR beside
DR. Reading it changes what two of the baselines look like.

\begin{table}[t]
\caption{\textbf{Headline comparison, multi-seed.} CICIoT2023, IID, $m{=}0.2$,
$p{=}0.1$; every entry is mean\,$\pm$\,s.d.\ over \textbf{5 seeds} (42--46).
Clean CA/macro-F1 are the no-attack models. \textbf{DR is reported with FPR
throughout}: FLAME's low BSR is bought by flagging two honest clients in five,
and on low-salience triggers GradCos flags honest clients slightly more often
than malicious ones. FLTrust never zeroes a client out (DR$=$FPR$=0$); its
robustness is reweighting, not exclusion. Best clean CA and best BSR per column
in \textbf{bold}.}
\label{tab:main}
\centering
\footnotesize
\begin{tabular}{@{}lcc|ccc|ccc@{}}
\toprule
 & Clean CA & Clean F1 & \multicolumn{3}{c|}{BSR by trigger} & \multicolumn{3}{c}{DR / FPR by trigger} \\
\cmidrule(lr){4-6}\cmidrule(l){7-9}
Defense & (no atk) & (no atk) & Fixed & Statistical & Temporal & Fixed & Statistical & Temporal \\
\midrule
FedAvg  & $\msCleanCAfedavg$ & $\msCleanFonefedavg$ & $\msBSRfedavgFixed$ & $\msBSRfedavgStatistical$ & $\msBSRfedavgTemporal$ & --- & --- & --- \\
Krum    & $\msCleanCAkrum$ & $\msCleanFonekrum$ & $\msBSRkrumFixed$ & $\msBSRkrumStatistical$ & $\msBSRkrumTemporal$ & $\msDRkrumFixed/\msFPRkrumFixed$ & $\msDRkrumStatistical/\msFPRkrumStatistical$ & $\msDRkrumTemporal/\msFPRkrumTemporal$ \\
FLAME   & $\msCleanCAflame$ & $\msCleanFoneflame$ & $\msBSRflameFixed$ & $\msBSRflameStatistical$ & $\msBSRflameTemporal$ & $\msDRflameFixed/\mathbf{\msFPRflameFixed}$ & $\msDRflameStatistical/\msFPRflameStatistical$ & $\msDRflameTemporal/\msFPRflameTemporal$ \\
FLTrust$^\dagger$ & $\msCleanCAfltrust$ & $\msCleanFonefltrust$ & $\mathbf{\msBSRfltrustFixed}$ & $\mathbf{\msBSRfltrustStatistical}$ & $\mathbf{\msBSRfltrustTemporal}$ & $\msDRfltrustFixed/\msFPRfltrustFixed$ & $\msDRfltrustStatistical/\msFPRfltrustStatistical$ & $\msDRfltrustTemporal/\msFPRfltrustTemporal$ \\
FoolsGold & $\msCleanCAfoolsgold$ & $\msCleanFonefoolsgold$ & $\msBSRfoolsgoldFixed$ & $\msBSRfoolsgoldStatistical$ & $\msBSRfoolsgoldTemporal$ & $\msDRfoolsgoldFixed/\msFPRfoolsgoldFixed$ & $\msDRfoolsgoldStatistical/\msFPRfoolsgoldStatistical$ & $\msDRfoolsgoldTemporal/\msFPRfoolsgoldTemporal$ \\
\textbf{GradCos} & $\mathbf{\msCleanCAgradcos}$ & $\mathbf{\msCleanFonegradcos}$ & $\mathbf{\msBSRgradcosFixed}$ & $\msBSRgradcosStatistical$ & $\msBSRgradcosTemporal$ & $\mathbf{\msDRgradcosFixed/\msFPRgradcosFixed}$ & $\msDRgradcosStatistical/\msFPRgradcosStatistical$ & $\msDRgradcosTemporal/\msFPRgradcosTemporal$ \\
\bottomrule
\end{tabular}
\\[3pt]
\raggedright\footnotesize $^\dagger$FLTrust is granted a clean server root set the
others do not have. Standard deviations over the five seeds are below $0.02$
except: FoolsGold fixed BSR $\pm0.489$, FLAME statistical BSR $\pm0.137$ and
temporal $\pm0.225$, FLTrust temporal $\pm0.147$, FLAME DR $\pm0.14$
(Sec.~\ref{sec:seed}).
\end{table}

\textbf{FLAME's robustness is a false-positive story.} It reaches BSR $\msBSRflameFixed$ on
the fixed trigger, the second-best number in the column, while flagging
$\msFPRflameFixed\pm\msFPRsdflameFixed$ of \emph{honest} clients---and
$\msFPRflameStatistical$--$\msFPRflameTemporal$ on the other two triggers, where
its DR is only $\msDRflameStatistical$--$\msDRflameTemporal$. A defense
discarding two honest updates in five is training on roughly $60\%$ of the
federation, which is a sufficient explanation for its $\msCleanCAflame$ clean
accuracy against FedAvg's $\msCleanCAfedavg$.
This is the mechanism behind the ``no free lunch'' observation, and it is
invisible without FPR: on BSR alone FLAME looks like a near-solution.

\textbf{GradCos's detector is not merely blind on low-salience triggers---it is
uninformative.} On statistical and temporal its DR ($\msDRgradcosStatistical$,
$\msDRgradcosTemporal$) is \emph{below} its FPR ($\msFPRgradcosStatistical$,
$\msFPRgradcosTemporal$). The clique test is not failing to find the
colluders so much as flagging clients at close to chance. We think this is the
sharpest available statement of the blind spot, and it is stronger evidence than
a low DR alone would be, because it rules out the reading that the detector is
simply conservative. Fig.~\ref{fig:drfpr} plots the two rates against each other
and makes the separation visible: only GradCos on the fixed trigger reaches the
top-left corner, while its low-salience points sit on the chance diagonal and
FLAME sits far to the right on every trigger.

\begin{figure}[t]
\centering
\includegraphics[width=0.68\linewidth]{fig_dr_fpr.pdf}
\caption{\textbf{Detection rate against false-positive rate} (CICIoT2023, IID,
$m{=}0.2$, $p{=}0.1$, 5 seeds). Only GradCos on the fixed trigger reaches the
top-left corner a detector should occupy. FLAME sits far right on every
trigger---its low backdoor rate is paid for by discarding honest updates---and
GradCos on the low-salience triggers sits \emph{on} the chance diagonal, flagging
honest and malicious clients at indistinguishable rates. FedAvg and the
coordinate-wise rules never flag a client and do not appear.}
\label{fig:drfpr}
\end{figure}

\textbf{FLTrust never excludes anyone.} Its DR and FPR are identically zero on
every trigger: no client's update ever receives a trust score of exactly zero.
Whatever robustness it has comes from continuous reweighting toward the server's
root direction, not from exclusion---which is also why it degrades the clean
model, since honest updates that disagree with a 300-sample root direction are
attenuated too.

\subsection{Salience decides the outcome, and it is measurable}
\label{sec:salience}

Grouped by trigger rather than by attacker count, the pattern in
Table~\ref{tab:main} is stark. On the \emph{fixed} trigger the direction- and
reference-based defenses work: GradCos $\msBSRgradcosFixed$ and FLTrust
$\msBSRfltrustFixed$ are lowest, FLAME $\msBSRflameFixed$, against FedAvg's
$\msBSRfedavgFixed$. On \emph{statistical} and \emph{temporal} they fail almost
completely---GradCos, Krum and FoolsGold all at $\msBSRgradcosStatistical$, FLAME
at $\msBSRflameStatistical$---with FLTrust on statistical
($\msBSRfltrustStatistical$) the lone exception, and it uses its clean root direction rather than inter-client
similarity.

The mechanism is visible in the update geometry we log. The fixed trigger stamps
header/flag features that few benign flows exercise, so learning it bends a
poisoned client's update in a direction honest clients never take. The
statistical and temporal triggers perturb volume, dispersion and timing features
and leave the poisoned update nearly collinear with an honest one, so the cosine
signal every one of these defenses depends on is simply not there. That FoolsGold---built
precisely to catch colluding sybils by cosine similarity---also fails
underlines that the blind spot belongs to the \emph{signal}, not to one defense's
tuning.

\textbf{Identifying the axis.} A hostile reading of the paragraph above is that
it is nearly definitional: ``direction-based defenses fail when the trigger does
not move the direction'' restates the mechanism rather than testing it. Two
things separate the claim from a tautology. First, $S$
(Eq.~\ref{eq:salience}) is measured independently of any defense's decision, so
BSR and DR can be regressed on it. Second, the three trigger families differ in
\emph{which} features they touch as well as how hard, so a two-group contrast
cannot separate feature coverage from magnitude. We therefore sweep the trigger
magnitude at a fixed feature set, $0.25\times$ to $3\times$ the base pattern, and
report BSR and DR against measured $S$ rather than against the trigger's name
% RETIRES WITH: lane 2, magnitude-sweep
(Fig.~\ref{fig:dose}). \PENDING

\begin{figure}[t]
\centering
% Self-upgrading: paper/make_figs.py writes figs/fig_dose.pdf once lane 2 of
% run_remediation.sh has produced at least three trigger magnitudes. Until then
% this falls back to the placeholder and the caption says so, so the switch
% cannot be forgotten.
\IfFileExists{figs/fig_dose.pdf}{%
  \includegraphics[width=0.95\linewidth]{fig_dose.pdf}%
  \caption{\textbf{Dose--response along the measured salience axis} (CICIoT2023,
  IID, $m{=}0.2$, $p{=}0.1$). BSR (left) and detection rate (right) against the
  salience gap $S$ the server actually observed, for trigger magnitudes
  $0.25$--$3\times$ at a fixed feature set. Marker shape is the trigger family
  ($\circ$ fixed, $\square$ statistical, $\triangle$ temporal); magnitude is swept
  within a family, so each curve is one identified axis rather than a contrast
  between families.}%
}{%
  \includegraphics[width=0.68\linewidth]{fig_detection.pdf}%
% RETIRES WITH: lane 2, magnitude-sweep
  \caption{\textbf{Dose--response along the measured salience axis.} \PENDING{}
  Placeholder shows the current detection-vs-BSR figure. The final version plots
  BSR and DR against the measured salience gap $S$ for trigger magnitudes
  $0.25$--$3\times$ at a fixed feature set, with the clean-model BSR overlaid to
  mark where a large trigger stops being a learned backdoor and becomes an
  out-of-distribution artifact.}%
}
\label{fig:dose}
\end{figure}

\subsection{The collusion signal is cheap to remove}
\label{sec:adaptive}

Every direction-based rule here detects colluders because colluders resemble one
another: GradCos's clique test flags a client for having a near-twin, FoolsGold
penalizes similar histories, FLAME's clustering separates a coherent minority.
That shared premise is the attack surface, and removing it costs the attacker
almost nothing.

Let $M$ be the malicious cohort and $\delta_j$, $j\in M$, their updates. The
attacker draws independent isotropic noise $\nu_j$ scaled to $\lambda\|\delta_j\|$
and subtracts the cohort mean,
\begin{equation}
\label{eq:evasion}
\tilde\delta_j \;=\; \delta_j + \Big(\nu_j - \tfrac{1}{|M|}\textstyle\sum_{k\in M}\nu_k\Big),
\end{equation}
so the perturbation sums to zero over $M$ and contributes nothing to the
aggregate, while each pair of malicious updates is decorrelated. There is no
optimization and one scalar knob, $\lambda$. We call this variant
\texttt{decorrelated}; it inflates the malicious update norms, so a magnitude
rule could still see it. The \texttt{adaptive} variant additionally rescales each
malicious update to the round's median honest norm, hiding that tell as well.

The evasion does what it is designed to do, and this can be checked without
running a single defense, because the quantity it targets is one we log. Measured
on the malicious block, the mean off-diagonal cosine among colluders falls from
$0.99$ to $0.16$--$0.33$ at $\lambda=1$ while the honest cohort's stays at
$0.99$; the median malicious-to-honest norm ratio stays near $1$ under the
norm-matched variant, against $\approx1.23$ for the noise-only variant, which is
the magnitude tell the norm-matching exists to hide; and the malicious mean
update---the component the backdoor rides on---is preserved. The colluders stop looking like each other and keep
contributing the same thing on average, which is precisely the property every
clique-, cluster- or history-similarity test is built to exclude.

\emph{Provenance.} These geometry figures come from a reduced-scale validation
run (6 clients, 3 rounds, 8k rows) used to verify the transform end to end. They
establish that it has the intended effect on the signal; they are \emph{not}
evidence about backdoor outcomes. Full-scale geometry, and every BSR/DR figure
% RETIRES WITH: lane 1a, adaptive-vs-direction (full-scale geometry)
below, come from the queued runs. \PENDING

Table~\ref{tab:adaptive} reports BSR, DR and FPR for the static and adaptive
attackers at the headline point against every direction-based defense, and
% RETIRES WITH: lanes 1a + 1b, adaptive-vs-direction and evasion-noise-sweep
Fig.~\ref{fig:evasion} sweeps $\lambda$. \PENDING

\begin{table}[t]
\caption{\textbf{Static vs.\ adaptive attacker}, CICIoT2023, IID, fixed trigger,
$m{=}0.2$, $p{=}0.1$. The adaptive attacker decorrelates the colluders
% RETIRES WITH: lane 1a, adaptive-vs-direction
(Eq.~\ref{eq:evasion}, $\lambda=1$) and norm-matches them. \PENDING}
\label{tab:adaptive}
\centering
\footnotesize
\begin{tabular}{@{}lcccc@{}}
\toprule
 & \multicolumn{2}{c}{static} & \multicolumn{2}{c}{adaptive} \\
\cmidrule(lr){2-3}\cmidrule(l){4-5}
Defense & BSR & DR/FPR & BSR & DR/FPR \\
\midrule
Krum      & $\msBSRkrumFixed$ & $\msDRkrumFixed/\msFPRkrumFixed$ & $\adBSRkrumAdaptive$ & $\adDRkrumAdaptive/\adFPRkrumAdaptive$ \\
FLAME     & $\msBSRflameFixed$ & $\msDRflameFixed/\msFPRflameFixed$ & $\adBSRflameAdaptive$ & $\adDRflameAdaptive/\adFPRflameAdaptive$ \\
FLTrust   & $\msBSRfltrustFixed$ & $\msDRfltrustFixed/\msFPRfltrustFixed$ & $\adBSRfltrustAdaptive$ & $\adDRfltrustAdaptive/\adFPRfltrustAdaptive$ \\
FoolsGold & $\msBSRfoolsgoldFixed$ & $\msDRfoolsgoldFixed/\msFPRfoolsgoldFixed$ & $\adBSRfoolsgoldAdaptive$ & $\adDRfoolsgoldAdaptive/\adFPRfoolsgoldAdaptive$ \\
GradCos   & $\msBSRgradcosFixed$ & $\msDRgradcosFixed/\msFPRgradcosFixed$ & $\adBSRgradcosAdaptive$ & $\adDRgradcosAdaptive/\adFPRgradcosAdaptive$ \\
GradCos-residual & $\resBSRSimpleFixed$ & $\resDRSimpleFixed/\resFPRSimpleFixed$ & $\resBSRAdaptiveFixed$ & $\resDRAdaptiveFixed/\resFPRAdaptiveFixed$ \\
\bottomrule
\end{tabular}
\end{table}

\begin{figure}[t]
\centering
% Self-upgrading, as with Fig.~\ref{fig:dose}: written by paper/make_figs.py once
% lane 1b of run_remediation.sh lands.
\IfFileExists{figs/fig_evasion.pdf}{%
  \includegraphics[width=0.95\linewidth]{fig_evasion.pdf}%
  \caption{\textbf{Evasion dose--response} (CICIoT2023, IID, fixed trigger,
  $m{=}0.2$, $p{=}0.1$). Backdoor success (left) against the decorrelation scale
  $\lambda$, with $\lambda{=}0$ the static attacker, and the collusion signal the
  defenses read (right): the mean pairwise cosine among colluders, which is the
  mechanism behind the left panel. The signal every clique-, cluster- and
  history-similarity rule depends on is gone well before $\lambda$ reaches the
  update's own scale.}%
}{%
  \includegraphics[width=0.88\linewidth]{fig_multiseed_headline.pdf}%
% RETIRES WITH: lane 1b, evasion-noise-sweep
  \caption{\textbf{Evasion dose--response.} \PENDING{} Placeholder. Final version:
  BSR and DR against the decorrelation scale $\lambda\in\{0,0.25,0.5,1,2\}$, with
  the measured malicious pairwise cosine on a second axis, showing how much noise
  the attacker needs before the clique signal disappears.}%
}
\label{fig:evasion}
\end{figure}

We expect, and will report either way, that the residual variant is \emph{also}
vulnerable: mean-centering removes the shared task direction but the decorrelation
noise survives it, so the residual clique signal should degrade for the same
reason the full-vector one does. If so, the honest conclusion is that our
diagnosis of the blind spot is correct and our proposed fix for it is not
adaptive-secure---which is a more useful thing to report than a defense evaluated
only against the attacker it was designed for.

\subsection{Do magnitude-based rules share the blind spot?}
\label{sec:magnitude}

If low-salience triggers evade defenses by not moving the update direction, then
rules that never read direction should not fail for that reason. This is the
control our earlier framing lacked: a claim that ``the whole cosine family has a
structural blind spot'' is only interesting against comparators outside the
family. Table~\ref{tab:magnitude} reports the coordinate-wise median, trimmed
mean, RFA, Bulyan and norm-clip at the headline point on all three triggers.
% RETIRES WITH: lane 3, magnitude-based-defenses
\PENDING

\begin{table}[t]
\caption{\textbf{Magnitude and order-statistic rules} at the headline point
(CICIoT2023, IID, $m{=}0.2$, $p{=}0.1$). None reads update direction. DR/FPR are
% RETIRES WITH: lane 3, magnitude-based-defenses
N/A for the coordinate-wise rules, which never single out a client. \PENDING}
\label{tab:magnitude}
\centering
\footnotesize
\begin{tabular}{@{}lccc@{}}
\toprule
Defense & Fixed & Statistical & Temporal \\
\midrule
Median        & $\mgBSRmedianFixed$ & $\mgBSRmedianStatistical$ & $\mgBSRmedianTemporal$ \\
Trimmed mean  & $\mgBSRtrimmedmeanFixed$ & $\mgBSRtrimmedmeanStatistical$ & $\mgBSRtrimmedmeanTemporal$ \\
RFA           & $\mgBSRrfaFixed$ & $\mgBSRrfaStatistical$ & $\mgBSRrfaTemporal$ \\
Bulyan        & $\mgBSRbulyanFixed$ & $\mgBSRbulyanStatistical$ & $\mgBSRbulyanTemporal$ \\
Norm-clip     & $\mgBSRnormclipFixed$ & $\mgBSRnormclipStatistical$ & $\mgBSRnormclipTemporal$ \\
\bottomrule
\end{tabular}
\end{table}

\subsection{Recovering detection: the residual variant}
\label{sec:residual}

If the blind spot is caused by the shared main-task gradient drowning the
backdoor component, mean-centering the updates before the cosine test should
expose it. It does---for detection. Table~\ref{tab:residual} compares the default
\texttt{clique} GradCos with the \texttt{residual} variant at the headline point
over five seeds. On both low-salience triggers the residual variant lifts the
malicious-client detection rate from $\msDRgradcosStatistical$ to
$\resDRSimpleStatistical$---and, on the standard we set in Sec.~\ref{sec:fpr}, it
does so honestly: its false-positive rate is
$\resFPRSimpleStatistical$--$\resFPRSimpleTemporal$, against
$\msFPRgradcosStatistical$ for the clique rule it replaces and
$\msFPRflameFixed$ for FLAME. Where the clique detector was at chance on these triggers, the
residual detector separates the colluders cleanly.

Detection does not translate into suppression equally. On the \emph{statistical}
trigger the recovered detection is enough: BSR falls from
$\msBSRgradcosStatistical$ to $\resBSRSimpleStatistical\pm\resBSRsdSimpleStatistical$.
On the \emph{temporal} trigger the backdoor is only partly contained
($\msBSRgradcosTemporal\rightarrow\resBSRSimpleTemporal\pm\resBSRsdSimpleTemporal$)
\emph{despite} detection at $\resDRSimpleTemporal$.

We previously reported this gap as an open question. It is a real hole in the
mechanistic story and we now say so directly: if the blind spot were fully
explained by ``the shared gradient masks the backdoor'', restoring detection to
$\resDRSimpleTemporal$ should restore suppression, and for temporal it does not. Two candidate
explanations are consistent with the per-round logs, and we can separate neither
with the evidence we have. First, at $20\%$ malicious the cohort mean used for
centering is itself contaminated, so mean-centering removes the backdoor component
unevenly. Second, filtering is per-round while the backdoor is persistent---durability
under continued training is the property Neurotoxin~\cite{zhang2022neurotoxin}
studies directly---so a handful of early misses is enough, and the
residual BSR does rise in \emph{late} rounds while DR stays near
$\resDRSimpleTemporal$, which is
consistent with early leakage compounding rather than a detection lapse. Either
way, per-round filtering against a persistent backdoor is structurally fragile,
and the residual variant should be read as evidence about the \emph{signal}
(collusion is recoverable once the common component is removed) rather than as a
working defense. Note also that the temporal trigger is the one our realizability
audit finds jointly implausible (Sec.~\ref{sec:realizable}), so this particular
gap may not correspond to any realizable attack.

\begin{table}[t]
\caption{GradCos \texttt{clique} vs.\ \texttt{residual} at the headline point
(CICIoT2023, IID, $m{=}0.2$, $p{=}0.1$; mean over 5 seeds). Mean-centering
restores detection on the low-salience triggers at a low false-positive rate;
suppression follows fully for statistical, partly for temporal. Mean over 5 seeds;
BSR s.d.\ is $\resBSRsdSimpleStatistical$ (statistical) and
$\resBSRsdSimpleTemporal$ (temporal).}
\label{tab:residual}
\centering
\footnotesize
\begin{tabular}{@{}lccc|ccc@{}}
\toprule
 & \multicolumn{3}{c|}{\texttt{clique} (default)} & \multicolumn{3}{c}{\texttt{residual}} \\
\cmidrule(lr){2-4}\cmidrule(l){5-7}
Trigger & BSR & DR & FPR & BSR & DR & FPR \\
\midrule
Fixed       & $\msBSRgradcosFixed$ & $\msDRgradcosFixed$ & $\msFPRgradcosFixed$
            & $\resBSRSimpleFixed$ & $\resDRSimpleFixed$ & $\resFPRSimpleFixed$ \\
Statistical & $\msBSRgradcosStatistical$ & $\msDRgradcosStatistical$ & $\msFPRgradcosStatistical$
            & $\mathbf{\resBSRSimpleStatistical}$ & $\mathbf{\resDRSimpleStatistical}$ & $\mathbf{\resFPRSimpleStatistical}$ \\
Temporal    & $\msBSRgradcosTemporal$ & $\msDRgradcosTemporal$ & $\msFPRgradcosTemporal$
            & $\resBSRSimpleTemporal$ & $\mathbf{\resDRSimpleTemporal}$ & $\mathbf{\resFPRSimpleTemporal}$ \\
\bottomrule
\end{tabular}
\end{table}

\subsection{Collusion: the malicious-ratio axis}
\label{sec:collusion}

Within the fixed trigger, sweeping the malicious fraction exposes how the
defenses degrade (Table~\ref{tab:ratio}). Krum is the cautionary case: robust at
$10\%$ ($\ratBSRkrumTen$) but collapsing to $\ratBSRkrumTwenty$ at $20\%$ and
$\ratBSRkrumThirty$ at $30\%$, because once the
malicious set is large enough its members become each other's nearest neighbours
and Krum selects one of them. GradCos, FLAME and FLTrust hold across the range.

This sweep is \textbf{single-seed}, and by our own variance argument
(Sec.~\ref{sec:seed}) single-seed numbers are least trustworthy exactly at
threshold-crossing points like Krum's collapse. We therefore treat it as
directional evidence for an ordering (Krum degrades with $m$ where the others do
not), not as a measurement of where the collapse happens; a three-seed repeat of
% RETIRES WITH: lane 5c, ratio-sweep-seeds
this sweep is in the released run plan. \PENDING

\begin{table}[t]
\caption{CICIoT2023, IID, fixed trigger: mean BSR vs.\ malicious fraction
(averaged over poison rates; \textbf{single-seed}, seed 42). Directional only---see
Sec.~\ref{sec:seed}.}
\label{tab:ratio}
\centering
\footnotesize
\begin{tabular}{@{}lccc@{}}
\toprule
Defense & $m=0.1$ & $m=0.2$ & $m=0.3$ \\
\midrule
FedAvg  & \ratBSRfedavgTen & \ratBSRfedavgTwenty & \ratBSRfedavgThirty \\
Krum    & \textbf{\ratBSRkrumTen} & \ratBSRkrumTwenty & \ratBSRkrumThirty \\
FLAME   & \ratBSRflameTen & \textbf{\ratBSRflameTwenty} & \ratBSRflameThirty \\
FLTrust & \textbf{\ratBSRfltrustTen} & \textbf{\ratBSRfltrustTwenty} & \textbf{\ratBSRfltrustThirty} \\
GradCos & \ratBSRgradcosTen & \ratBSRgradcosTwenty & \ratBSRgradcosThirty \\
\bottomrule
\end{tabular}
\end{table}

\subsection{No free lunch: the clean-accuracy cost}
\label{sec:tradeoff}

Fig.~\ref{fig:tradeoff} plots the two axes against each other over the breadth
grid; reading Table~\ref{tab:main} down the clean-CA column reframes the fixed-trigger
result. FLAME and FLTrust reach a low BSR with the weakest clean models
($\msCleanCAflame$ and $\msCleanCAfltrust$, macro-F1 $\msCleanFoneflame$ and
$\msCleanFonefltrust$) against FedAvg's $\msCleanCAfedavg/\msCleanFonefedavg$. GradCos, Krum and FoolsGold keep FedAvg-level clean accuracy. On
the fixed trigger GradCos ties FLTrust's BSR at $+\msCAadvGradcos$ clean accuracy.

Sec.~\ref{sec:fpr} supplies the mechanism for half of this: FLAME's deficit is
attributable to discarding a fraction $\msFPRflameFixed$ of honest updates, not
to its noise scale alone. That distinction matters for the criticism that our clean-accuracy
critique might be a tuning artifact---a badly chosen DP noise scale can crater
clean accuracy by itself. The FPR evidence shows the loss is happening at the
\emph{clustering} stage, upstream of the noise, and the calibration protocol of
Sec.~\ref{sec:tuning} tests the remaining exposure by sweeping the noise scale on
held-out data. Until that run lands we mark this claim as supported for the
% RETIRES WITH: the calibration run (FLAME noise-scale sweep)
clustering mechanism and provisional as to the noise scale. \PENDING

\begin{figure}[t]
\centering
\includegraphics[width=0.68\linewidth]{fig_tradeoff.pdf}
\caption{Clean accuracy vs.\ fixed-trigger backdoor success (CICIoT2023, IID;
single-seed breadth view). Bottom-right is best.}
\label{fig:tradeoff}
\end{figure}

\subsection{FLNET2023, and where GradCos is unreliable}

FLNET2023 is an easier task---clean accuracy $\flCleanCAgradcos$ for FedAvg,
Krum and GradCos---which sharpens the clean-accuracy point: FLAME's and FLTrust's
clean macro-F1 fall to $\flCleanFoneflame$ and $\flCleanFonefltrust$ on a
near-separable dataset (Table~\ref{tab:flnet}). At the headline fixed-trigger
point FedAvg is fully backdoored ($\flBSRfedavgFixed$) while Krum
($\flBSRkrumFixed$), FLAME ($\flBSRflameFixed$) and FLTrust
($\flBSRfltrustFixed$) suppress it.

GradCos does not, on the worst-case seed, and we want to state the consequence
rather than bury it in a footnote. The FLNET GradCos fixed $m0.2p0.1$ cell gives
BSR $\{0.98, 0.18, 0.008, 0.008\}$ across seeds. That is the same
locks-on-or-does-not bimodality we use to disqualify FoolsGold on CICIoT
(Sec.~\ref{sec:seed}). \textbf{So ``GradCos reliable, FoolsGold unreliable'' is a
dataset-dependent statement, not a property of the method}: it holds on CICIoT
(s.d.\ $\msBSRsdgradcosFixed$) and fails on FLNET. Both methods depend on a clique separating
early enough, and on a near-separable dataset where honest updates agree strongly,
whether it separates is close to a coin flip. We retain the worst-case draw in
the table.

\begin{table}[t]
\caption{FLNET2023, IID (single-seed). FLAME/FLTrust keep a much weaker clean
model on a near-separable task; GradCos is bimodal across seeds here.}
\label{tab:flnet}
\centering
\footnotesize
\begin{tabular}{@{}lccc@{}}
\toprule
Defense & Clean CA & Clean F1 & BSR (fixed, $m0.2p0.1$) \\
\midrule
FedAvg  & \flCleanCAfedavg & \flCleanFonefedavg & \flBSRfedavgFixed \\
Krum    & \flCleanCAkrum & \flCleanFonekrum & $\mathbf{\flBSRkrumFixed}$ \\
FLAME   & \flCleanCAflame & \flCleanFoneflame & \flBSRflameFixed \\
FLTrust & \flCleanCAfltrust & \flCleanFonefltrust & \flBSRfltrustFixed \\
GradCos & $\mathbf{\flCleanCAgradcos}$ & $\mathbf{\flCleanFonegradcos}$ & \flBSRgradcosFixed$^{\ast}$ \\
\bottomrule
\end{tabular}
\\[2pt]
\raggedright\footnotesize $^{\ast}$Worst-case of $\{0.98,0.18,0.008,0.008\}$
across seeds---bimodal, see text.
\end{table}

\subsection{Non-IID}

At the single non-IID (Dirichlet, $\alpha=0.5$) configuration we test---fixed
trigger, $m=0.2,p=0.1$ (Table~\ref{tab:noniid})---GradCos is the strongest
defense: BSR
$\niBSRgradcosFixed$ at clean accuracy $\niCleanCAgradcos$. FedAvg is backdoored
($\niBSRfedavgFixed$) and so is Krum ($\niBSRkrumFixed$):
heterogeneous honest data already spreads updates out, which blunts Krum's
distance test.

We previously wrote that this shows ``the mechanism survives heterogeneity.'' It
does not show that. One Dirichlet concentration, one seed, one trigger and one
operating point cannot characterize behavior in the setting that matters most
for a federated IDS, and the one thing this slice does establish---that
heterogeneity flips Krum from robust to backdoored---is itself evidence that
non-IID data changes the geometry these defenses read. The central claim of this
paper is therefore \emph{untested} under realistic heterogeneity, and we say so.
A sweep over $\alpha\in\{0.1,0.5,1.0\}$ across all three triggers is in the
% RETIRES WITH: lane 6, non-IID breadth
released run plan. \PENDING

\begin{table}[t]
\caption{CICIoT2023, non-IID (Dirichlet $\alpha=0.5$), fixed trigger,
$m=0.2,p=0.1$ (\textbf{single-seed, single configuration}; not a
characterization of non-IID behavior---see text).}
\label{tab:noniid}
\centering
\footnotesize
\begin{tabular}{@{}lccc@{}}
\toprule
Defense & Clean CA & Attacked CA & BSR \\
\midrule
FedAvg  & \niCleanCAfedavg & \niCAfedavgFixed & \niBSRfedavgFixed \\
Krum    & \niCleanCAkrum & \niCAkrumFixed & \niBSRkrumFixed \\
FLAME   & \niCleanCAflame & \niCAflameFixed & \niBSRflameFixed \\
FLTrust & \niCleanCAfltrust & \niCAfltrustFixed & \niBSRfltrustFixed \\
\textbf{GradCos} & $\mathbf{\niCleanCAgradcos}$ & $\mathbf{\niCAgradcosFixed}$ & $\mathbf{\niBSRgradcosFixed}$ \\
\bottomrule
\end{tabular}
\end{table}

\subsection{Reproducibility and variance}
\label{sec:seed}

The headline comparison is run over five seeds and reported with dispersion;
the breadth grid, the ratio sweep, and the FLNET and non-IID slices remain
single-seed and are labeled as such. Variance is not uniform, and where it is
large it is itself a finding. \emph{FoolsGold on the fixed trigger} has s.d.
$\msBSRsdfoolsgoldFixed$ (BSR range
$\msBSRminfoolsgoldFixed$--$\msBSRmaxfoolsgoldFixed$): it either locks onto the colluders' shared
history early or never does. \emph{GradCos on FLNET} behaves the same way. By
contrast GradCos and FLTrust on CICIoT fixed, and every defense on the
low-salience triggers, have s.d.\ $\leq0.02$---so the salience finding is not a
seed artifact.

There is a tension here we should not paper over. We use the FoolsGold result to
argue that single-seed numbers are untrustworthy at threshold-crossing operating
points, and most of our breadth evidence is single-seed. The two are consistent
only under a distinction we now make explicitly: dispersion is large at
\emph{borderline} points, where the outcome hinges on whether a detector locks on
in time, and small where BSR is pinned near $0$ or $1$. Our multi-seed evidence
covers the borderline cases that matter; the single-seed grid supplies breadth,
and we draw only ordering claims from it. Where a single-seed cell does drive a
conclusion---Krum's collapse point, the FLNET comparison, the non-IID result---we
now say so in the caption and mark the multi-seed repeat as outstanding.

\section{Discussion and Limitations}

\textbf{The blind spot is shared, and it is the point.} GradCos is not a
universal winner and we do not present it as one. Its practical claim is narrow:
it removes FLTrust's clean-data assumption and Krum's need to know the attacker
count, it holds where Krum collapses, and it keeps the clean accuracy FLAME and
FLTrust give up. On low-salience triggers it fails, and so do Krum, FLAME and
FoolsGold, for the same reason---they all read update direction. That common
failure is the more important result, because it says the family has a structural
blind spot exactly where a competent attacker would aim.

\textbf{And the blind spot is not the worst of it.} A defense that fails against
a trigger it cannot see is a bounded problem; a defense whose detection signal an
attacker can delete for free is a different kind of problem.
Sec.~\ref{sec:adaptive} deletes it with mean-zero noise and one scalar. We
consider this the strongest reason to be skeptical of clique-, cluster- and
history-similarity detectors in FL, including our own, and it is why we frame
GradCos's contribution as diagnostic. The residual variant shows the collusion
signal is recoverable once the shared task component is removed---which is a real
finding about the signal---but recovery of a signal an attacker can decorrelate
is not a defense.

\textbf{What would actually help.} Three directions follow from our own negative
results rather than from speculation. Aggregating evidence \emph{across} rounds
rather than filtering per round addresses the persistence asymmetry in
Sec.~\ref{sec:residual}: the backdoor is cumulative and the detector is not.
Pairing a direction test with a magnitude or coordinate test attacks the
decorrelation evasion from the side it does not cover, since zero-sum noise
inflates norms unless it is also norm-matched, and norm-matching perturbs the
aggregate. And any defense in this family should be reported with FPR, because
Sec.~\ref{sec:fpr} shows two established methods behave very differently once it
is visible.

\textbf{Scope, stated plainly.} Our evidence base is narrower than the generality
of the claims would ideally require, and the limitations compound rather than
sit independently:

\emph{Attack variants.} We sweep the simple data-poisoning variant at scale plus
the adaptive evasion of Sec.~\ref{sec:adaptive}. Norm-matching alone and the
distributed (DBA) variant are implemented and runnable but not swept, so our
static-attack claims are about direct poisoning.

\emph{Threat model.} Feature-space only. Sec.~\ref{sec:realizable} quantifies the
distance to realizable traffic rather than closing it, and finds our temporal
trigger jointly implausible by two orders of magnitude; conclusions resting on
that trigger are correspondingly weaker.

\emph{Base detector.} The CICIoT2023 grouped-task clean accuracy
($\approx0.70$, macro-F1 $\approx0.55$) reflects a deliberately small
$20$-client, $100$-round, $400$k-row budget. This is a genuine threat to
external validity and not merely a caveat about absolute numbers: update-direction
geometry in an underfit model may differ from a well-fitted detector's, and our
central claim is about that geometry. FLNET2023, where the same architecture
reaches macro-F1 $\flCleanFonegradcos$, is the partial control---the salience ordering holds
there too---but FLNET is near-separable and easier in a way that may not
generalize either. Whether the blind spot survives at, say, macro-F1 $0.85$ on a
hard task is open.

\emph{Heterogeneity.} One Dirichlet configuration; see the non-IID subsection.

\emph{Breadth.} One architecture, two datasets, one trigger family at a time.

\emph{Taxonomy.} Our direction-vs-magnitude split is a useful cut, not an
exhaustive one. FLDetector~\cite{zhang2022fldetector} scores clients on the
consistency of their updates \emph{across} rounds and so belongs to neither
family as we define them; our blind-spot claim is about the two families we
evaluate, not about every conceivable signal.

\emph{Why not certification?} Certified defenses such as CRFL~\cite{xie2021crfl}
bound backdoor success under a magnitude constraint on the adversary's
perturbation. That is exactly the regime our low-salience triggers are
constructed to sit inside, so a certificate of the CRFL form would be satisfied
by triggers that nevertheless reach BSR $1.0$ here. Certification bounds the
perturbation; our finding is that a small perturbation is enough.

\emph{Hyperparameters.} Fixed by the declared held-out protocol of
Sec.~\ref{sec:tuning}; the calibration run is outstanding, and until it lands the
grid uses pre-registered defaults.

\section{Conclusion}

We studied backdoors against a federated CNN\,+\,LSTM intrusion detector and the
Byzantine-robust defenses meant to stop them, across two 2023 benchmarks. The
factor that decides whether a similarity-based defense survives is how far the
trigger rotates the update off the honest consensus direction---measured, not
inferred from the trigger's name---and not the number of attackers. That signal
is also cheap for an attacker to remove: zero-sum noise across the colluding
cohort decorrelates them, preserves the aggregate, and defeats every
direction-based rule we test including our own \textsc{GradCos}, which otherwise
matches the strongest baseline on conspicuous triggers without a clean server
dataset and at $+\msCAadvGradcos$ clean accuracy. Reporting false-positive rates alongside
detection rates changes what two established baselines look like, and a
joint-plausibility audit shows the ``in-distribution trigger'' premise this
literature relies on does not follow from the marginal checks usually offered for
it. We release the pipeline, the per-round update-geometry logs and the audit, and
we regard the open problem the study surfaces---aggregation that is robust to
low-salience triggers \emph{and} to an attacker who knows what the aggregator
measures---as the useful output of the work.

%---------------------------------------------------------------------
%---------------------------------------------------------------------
%  Elsevier declarations. JISA requires each of these; the competing
%  interest declaration is additionally uploaded as a .docx through
%  Editorial Manager's declaration tool.
%---------------------------------------------------------------------

\section*{CRediT authorship contribution statement}

\textbf{Usama Mustafa:} Conceptualization, Methodology, Software,
Validation, Formal analysis, Investigation, Data curation, Writing --
original draft, Visualization.
\textbf{Imran Rashid:} Conceptualization, Methodology, Supervision,
Writing -- review and editing, Project administration.

\section*{Declaration of competing interest}

The authors declare that they have no known competing financial
interests or personal relationships that could have appeared to
influence the work reported in this paper.

\section*{Data availability}

Both datasets are public. CICIoT2023 is distributed by the Canadian
Institute for Cybersecurity~\cite{neto2023ciciot}; FLNET2023 is
distributed by the authors of~\cite{flnet2023}. The training pipeline,
the ten aggregation rules, the adaptive attacker, the trigger
realizability audit, the per-round update-geometry logs, and the result
files behind every number in this paper are available at
\url{\repourl}. Every table and every figure is regenerated from those
result files by \texttt{paper/make\_numbers.py} and
\texttt{paper/make\_figs.py}; no number in this paper is transcribed by
hand.

\section*{Funding}

This research did not receive any specific grant from funding agencies
in the public, commercial or not-for-profit sectors.


%---------------------------------------------------------------------
\bibliographystyle{elsarticle-num}
\bibliography{refs}

\end{document}
